% concatenated from 3-dim_PqN-v25.tex
 \documentclass[12pt]{article}
%\documentclass{article}
%\usepackage[all]{xy}
\usepackage{amssymb}
%\usepackage{amsfont}
\usepackage{amsmath}
\usepackage{amsthm}
\usepackage{calc}
\usepackage{stackrel}
%\usepackage{showkeys}
%\usepackage[italian]{babel}
%\usepackage[all]{xypic}
\usepackage{color}
\usepackage{cancel}
\usepackage{lineno}
\usepackage{hyperref}
%\linenumbers
\textwidth18.5cm
\topmargin-15mm
\evensidemargin-20mm
\oddsidemargin-10mm
 \textheight23.5cm

\usepackage[affil-it]{authblk}
%\usepackage[utf8]{inputenc}
 
% % % \textwidth17.5cm
% % % \textheight26cm
% % % \topmargin-25mm
% % % %\voffset-20mm
% % % \evensidemargin-20mm
% % % \oddsidemargin-15mm

\usepackage{graphicx}

\usepackage{comment} % package to comment between \begin{comment}...\end{comment}

\newtheorem{theorem}{Theorem}%[section]
\newtheorem{prop}[theorem]{Proposition}
\newtheorem{lem}[theorem]{Lemma}
\newtheorem{cor}[theorem]{Corollary}
\newtheorem{defi}[theorem]{Definition}
\newtheorem{remark}[theorem]{Remark}
\newenvironment{rem}{\begin{remark}\rm}{$\diamondsuit$\end{remark}}
\newtheorem{example}[theorem]{Example}
\newtheorem{conj}[theorem]{Conjecture}
\newenvironment{exa}{\begin{example}\rm}{$\diamondsuit$\end{example}}

\newtheorem{theoremblue}[theorem]{\textcolor{blue}{Theorem}}
\newtheorem{propblue}[theorem]{\textcolor{blue}{Proposition}}


%%%%%%%%%%%%%%%%%%%%%%%%%%%%%%%%%%%%%%%%%%%%%%%%%%%%%%%%%%%
\newcommand{\iu}{{\rm i}}
\newcommand{\imag}{{\rm i}}
\newcommand{\sech}{{\rm sech}}
\newcommand{\tepsilon}{\tilde{\epsilon}}
\renewcommand{\Re}{{\rm Re}}
\renewcommand{\Im}{{\rm Im}}
\newcommand{\sgn}{{\rm sgn}}
\newcommand{\cosec}{{\rm cosec}}
\newcommand{\eps}{{\epsilon}}
\newcommand{\req}[1]{(\ref{#1})}
%\newcommand{\bm}[1]{\mbox{\boldmath{$#1$}}}
%%%%%%%%%%%%%%%%%%%%%%%%%%%%%%%%%%%%%%%%%%%%%%%%%%%%%%%%%%%

%%%%%%%%%%%%%%%%%%%%%%%%%%%%%%%%%%%%%%%%%%%%%%%%%%%%%%%%%%%
\newcommand{\ua}{{\overline u}}
\newcommand{\rhob}{{\overline \rho}}
\newcommand{\pvi}{-\mskip -20mu \int}
\newcommand{\oper}{{\cal H}}
\newcommand{\ep}{\epsilon}
\newcommand{\drho}{\rho_{{}_\Delta}}
%\newcommand{\dPre}{{P_{{}_\Delta}}}
\newcommand{\dpre}{{P_{{}_\Delta}}}
\newcommand{\dbarpre}{{{\langle p\rangle}_{{}_\Delta}}}
\newcommand{\so}{{(0)}}
\newcommand{\si}{{(1)}}
\newcommand{\sii}{{(2)}}
%%%%%%%%%%%%%%%%%%%%%%%%%%%%%%%%%%%%%%%%%%%%%%%%%%%%%%%%%%%
\newcommand{\DP}[1]{{{\bar{P}_{(2)}\Big\vert_{-#1}^{+#1}}}}
%%%%%%%%%%%%%%%%%%%%%%%%%%%%%%%%%%%%%%%%%%%%%%%%%%%%%%%%%%%
\newcommand{\todo}[1]{
\vspace{5 mm}\par \noindent 
\framebox{
\begin{minipage}[c]{0.9 \textwidth} \small \tt #1
\flushright ...ToDo
\end{minipage}}
\vspace{5 mm}\par
}

%%%%%%%%%%%%%%%%%%%%%%%%%%%%%%%%%%%%%%%%%%%%%%%%%%%%%%%%%%%%%%%%%%%%%%

\def\d{{\mathrm d}}
\def\pam{{\beta}}
\def\lwp{{\epsilon}}
\def\ri{{\rm i}}
\def\rd{{\rm d}}
\def\upi{{\pi }}
\def\zl{{\zeta}}
\def\ztu{{{\eta}_1}}
\def\ztl{{{\eta}_2}}
\def\zu{{h}}
\def\vp{{\phi}}
\def\xx{{\xi}}
% \def\ou{{{\overline u}}}
% \def\ouu{{{\overline u_1}}}
% \def\ouj{{{\overline u_j}}}
% \def\oul{{{\overline u_2}}}
%\def\uu{{u_1}}
%\def\uuj{{u_j}}
%\def\ul{{u_2}}
%\def\pdoth{{\dot{\Pi_1}}}

\def\barr{\hbox{I\hskip -.5ex R}}
\def\sbarr{\hbox{{\tiny I}\hskip -.5ex {\tiny R}}}
%%%% Macros for Sect2-1%%%%%%%%%%%%%%%%%%%%%%%%%%%%%%%%%%%%%%%%%%%%%%%%%%%%%%%%%%%%%%%%%%
\usepackage{cancel}
\newcommand{\prvi}{-\mskip -20mu \int}
%\newcommand{\oper}{{\cal H}}
\newcommand{\mb}[1]{{\mathbf{#1}}}
\newcommand{\bs}[1]{{\boldsymbol{#1}}}
\newcommand{\bsl}[1]{{\boldsymbol {#1}}}
\def\parpo#1#2{\{#1,#2\}}
\def\ddd#1#2{\displaystyle{\frac{\partial #1}{\partial #2}}}
\def\fddd#1#2{\displaystyle{\frac{\delta #1}{\delta #2}}}
\newcommand{\DT}{\frac{d}{dt}}
\def\la{\lambda}        \def\La{\Lambda}
\newcommand{\dsl}[1]{{\displaystyle{#1}}}
\newcommand{\RR}{{{\mathbb{R}}}}
\newcommand{\brke}[2]{{{\langle #1,#2\rangle}}}

\def\M{{\mathcal M}}
\def\Tr{\operatorname{Tr}}
\newcommand{\rosso}[1]{\textcolor{red}{#1}}
\newcommand{\blu}[1]{\textcolor{blue}{#1}}

\begin{document}
\title{On three-dimensional Poisson quasi-Nijenhuis\\ manifolds and Haantjes structures}
%\title{Some remarks on Poisson quasi-Nijenhuis\\manifolds and Haantjes structures}
\date{} 
%\author{The authors}

\author{E.\ Chu\~no Vizarreta${}^{1}$, I.\ Mencattini${}^2$,  
M.\ Pedroni${}^{3,4}$}

\affil{
{\small  $^1$Unidade Acad\^emica de Belo Jardim, Universidade Federal Rural de Pernambuco, Brazil}\\
{\small eber.vizarreta@ufrpe.br  %(ORCID 0000-0002-4893-9186)
}\\
\medskip
{\small $^2$Instituto de Ci\^encias Matem\'aticas e de Computa\c c\~ao,  Universidade de S\~ao Paulo, Brazil}\\
{\small igorre@icmc.usp.br %(ORCID 0000-0002-1295-9396)
}\\
\medskip
{\small $^3$Dipartimento di Ingegneria Gestionale, dell'Informazione e della Produzione,  Universit\`a di Bergamo, Italy}\\
{\small marco.pedroni@unibg.it %(corresponding author, ORCID 0000-0002-7358-0945)}
}\\
\medskip
{\small  $^4$INFN, Sezione di Milano-Bicocca, Piazza della Scienza 3, 20126 Milano, Italy}
}

\maketitle
\abstract{\noindent
\begin{comment} %%%%%%%% Abstract as in v13
In this note we first characterize Poisson quasi-Nijenhuis structures on three-dimensional oriented manifolds %, 
whose underlying Poisson tensor never vanishes, % is maximally non-degenerate, 
and we show that each of these structures is involutive. Then we prove that every such three-dimensional 
Poisson quasi-Nijenhuis manifold is a Haantjes manifold 
and that %, under a further mild assumption on the underlying Poisson tensor, 
it carries a generalized Lenard-Magri chain.
\end{comment} 
In this note we first characterize Poisson quasi-Nijenhuis structures on three-dimensional oriented manifolds 
whose underlying Poisson tensor never vanishes. We then apply this result to show that each of these structures is  (locally) a deformation of a PN structure and is involutive. Finally, we prove that every such three-dimensional 
Poisson quasi-Nijenhuis manifold is a Haantjes manifold and that it carries a generalized Lenard-Magri chain.


\medskip\par\noindent
{\bf Keywords:} Integrable systems; Poisson quasi-Nijenhuis manifolds; Generalized Lenard-Magri chains; Haantjes tensors; Haantjes manifolds.
\medskip\par\noindent
{\bf MSC codes:} 37J35, 53D17, 70H06.}  

\baselineskip=0,6cm

\section{Introduction}

The notion of Poisson-Nijenhuis structure, i.e., a \emph{compatible} pair of a Poisson bivector $\pi$ and a 
torsionless tensor field $N$ of type $(1,1)$, is instrumental to the theory of classical integrable systems, since it provides a geometrical framework to produce involutive  families of conserved quantities for the underlying Hamiltonian dynamical system, see \cite{KM, Magri-Morosi}. It is important to stress that this property hinges crucially on the hypothesis that the Nijenhuis torsion of $N$ is equal to zero, see \eqref{eq:torsionA} and Section \ref{sec:2} for the relevant definitions. In fact, even a controlled change on this assumption destroys, in general, the possibility to produce the desired families of first integrals, spoiling a possible use of the newborn geometrical structure in the theory of integrable models.

An example of this state of affairs is provided by the so called Poisson quasi-Nijenhuis structures, i.e., a \emph{compatible} triple of a Poisson bivector $\pi$, a 3-form $\phi$ and a $(1,1)$ tensor field $N$, whose Nijenhuis torsion is given in  \eqref{eq:PqNtor}. %, see Section \ref{sec:2}. 
Although this class of structures represents a very interesting generalization of the Poisson-Nijenhuis one (see \cite{SX, Antunes2008, C-NdC-2010, BursztynDrummondNetto2021, DMP2024}) and has interesting applications in mathematical physics 
(see, for example, \cite{Zucchini}), it did not play any role in the theory of integrable systems until the discover 
of its relation with 
$n$-particle Toda systems, see \cite{FMOP2020,FMP2024}. In these papers the notion of \emph{involutive} Poisson quasi-Nijenhuis structure was introduced, see Definition \ref{def:invpqn} hereafter, and a set of \emph{sufficient} conditions for a Poisson quasi-Nijenhuis structure to be involutive were presented, see Theorem \ref{theo:inv} below. It is worth mentioning that, since these conditions are very stringent, the task of finding (interesting) examples of involutive Poisson quasi-Nijenhuis structures is a non-trivial open problem, see for example \cite{FMP2023}.

%Stealing 
In line with the title of \cite{Kos-Recursion}, the attempt to go beyond recursion operators to find a more general theoretical framework than the Poisson-Nijenhuis one %but still allowing the definition of chains of conserved quantities, 
was recently undertaken in \cite{MagriHaa, MagriVeselov, MagriSymm} and \cite{tempestatondoben, tempestatondo, tempestatondohigher,tempestatondoclass}. In these works, one of the main novelties is the presence of (families of) $(1,1)$ tensor fields whose \emph{Haantjes} torsion is equal to zero. 
As the Nijenhuis one, the Haantjes torsion of  $N$ is a 2-form with values in the tangent bundle of the underlying manifold, see  \eqref{eq:haant}, which controls the Frobenius integrability of the 2-planes generated by pairs of eigenvectors of $N$. 
Its relevance in the theory of integrable systems of hydrodynamic type was already observed in 
\cite{Bogo96-182,FeraMar}, but its importance in the theory of finite-dimensional integrable system was unnoticed until the above mentioned works of Magri, Tempesta and Tondo.

The present note is a \emph{first attempt} to construct a bridge between Poisson quasi-Nijenhuis and Haantjes structures, aiming to find, in a near future, more applications of the former to the theory of finite-dimensional integrable systems.
Our results concern oriented three-dimensional manifolds and generalize the one proved in \cite{Cheng-Sheng} for Poisson-Nijenhuis (PN from now on) manifolds. Assuming that the Poisson tensor never vanishes, we describe any Poisson quasi-Nijenhuis manifold in terms of few geometrical objects. Moreover, we show that such manifolds are all involutive. Finally, we show that they are Haantjes manifolds with a generalized Lenard-Magri chain.

A more detailed plan of the paper is as follows. In Section \ref{sec:2} we first collect some general information about Poisson quasi-Nijenhuis (PqN from now on) structures and we remind the definition of involutive PqN manifold. Then we recall two results, the first giving a set of sufficient conditions %guaranteeing 
for the involutivity of a PqN  structure, see Theorem \ref{theo:inv}, the second enclosing a general statement about 
deformations of PqN structures, see Theorem \ref{thm:defthe}. 
Section \ref{sec:3} is devoted to characterize %\textcolor{red}{(maximally non-degenerate)} 
PqN structures (with never vanishing Poisson tensor) on an oriented three-dimensional manifold, see %Proposition
Theorem \ref{prop 3D PqN}. %, and to prove that such structures are involutive, see %Theorem
As an application of this result we prove that such structures are (locally) obtained by deforming a PN structure via a closed 2-form, see Proposition \ref{prop:local}, and that they are involutive, see Proposition \ref{thm:inv3}. Moreover, in Section \ref{sec:4} they are shown to be Haantjes structures, see Theorem \ref{theo:haan}, and that they carry a generalized Lenard-Magri chain, see Theorem \ref{theo:Len-Ma-Ge}. In Subsection \ref{ss:tempton} we make a few final comments aiming to link the structures introduced in the recent works \cite{tempestatondoben, tempestatondo, tempestatondohigher,tempestatondoclass,TondoLag} with the geometry of the PqN 
manifolds studied in this note. 
In Section \ref{sec:Appendix} we collect the proof of an identity used in the main text, an alternative proof of Proposition \ref{thm:inv3} and few facts about the Haantjes torsion. We close this manuscript with Section \ref{sec:comfin}, which encloses a few final remarks about the results obtained in this paper and some possible future lines of investigation opened in this work.

\subsection{Notations and conventions} Hereafter we will systematically use the following conventions. All manifolds, 
vector bundles and their sections will be smooth.  The space of sections of the vector bundle $V$ will be denoted with $\Gamma(V)$, with the exception of the sections of $TM$, i.e., vector fields, and of $\Lambda^k(T^\ast M)$, i.e., $k$-differential forms, which will be denoted by $\mathfrak X(M)$ and, respectively, $\Omega^k(M)$. If $V$ is a vector bundle over a manifold $M$ and $A\in\text{End}_{C^\infty(M)}(V)$, $A^\ast$ will denote the endomorphism of $V^\ast$, the dual bundle of $V$, defined by $\langle A^\ast\eta,s\rangle=\langle\eta, A s\rangle$ for all $s\in\Gamma(V)$ and $\eta\in\Gamma(V^\ast)$. 
Every bivector field $P\in\Gamma(\Lambda^2(TM))$ defines a unique (and it is defined uniquely by) $P^\sharp\in\text{Hom}_{C^\infty(M)}(T^\ast M,TM)$ via the  
formula $\langle\beta,P^\sharp\alpha\rangle=P(\alpha,\beta)$, for all $\alpha,\beta\in\Omega^1(M)$. Similarly, every 2-form $\omega\in\Omega^2(M)$ defines a unique (and it is defined uniquely by) $\omega^\flat\in \text{Hom}_{C^\infty(M)}(TM,T^\ast M)$ via $\langle\omega^\flat X,Y\rangle=\omega(X,Y)$, for all $X,Y\in\mathfrak X(M)$. Finally, recall that to every $N\in\text{End}_{C^\infty(M)}TM$, i.e., a tensor field of type $(1,1)$, one can associate its Nijenhuis torsion $T_N$, which is the 2-form on $M$ with values in $TM$ defined by 
\begin{equation}
T_N(X,Y)=[NX,NY]-N([NX,Y]+[X,NY]-N[X,Y]).\label{eq:torsionA}
\end{equation}
The (1,1) tensor field $N$ will be called Nijenhuis torsion-free, or, more simply, torsion-free (or a Nijenhuis operator) 
if $T_N\equiv 0$.

\par\bigskip\noindent
{\bf Acknowledgments.} 
%We thank ???Murilo do Nascimento Luiz, Franco Magri, Giovanni Ortenzi, and Giorgio Tondo??? for useful discussions. %MP thanks the ICMC-UPS, {\em Instituto de Ci\^encias Matem\'aticas e de Computa\c c\~ao\/} of the University of S\~ao Paulo, and the Department of Mathematics and its Applications of the University of Milano-Bicocca for their hospitality, the Funda\c c\~ao de Amparo \`a Pesquisa do Estado de S\~ao Paulo - FAPESP, Brazil, for supporting his visit in 2023 to the ICMC-USP with the grant 2022/02454-8, and the University of Bergamo, for supporting his visit in 2024 to the ICMC-USP within the program {\em Outgoing Visiting Professors}.
%This project has received funding from the European Union's Horizon 2020 research and innovation programme under the Marie Sk{\l}odowska-Curie grant no 778010 {\em IPaDEGAN} as well as by the Italian PRIN 2022 (2022TEB52W) - PE1 - project {\em The charm of integrability: from nonlinear waves to random matrices}. 
All authors gratefully acknowledge the auspices of the GNFM Section of INdAM and of the Brazilian Centre of Geometry grant No. 2024/00923-6, Fapesp.
%, and the anonymous referee for suggesting some improvements in the presentation.
%We are grateful to the anonymous referee, whose suggestions helped us to substantially improve the content and the presentation of our manuscript.


\section{Recollection of Poisson quasi-Nijenhuis structures}\label{sec:2}

A Poisson-Nijenhuis structure on a manifold $M$ is a pair $(\pi,N)$ of a Poisson bivector  $\pi$ and a torsion-free tensor field $N$ of type $(1,1)$, satisfying the following compatibility conditions: 
\begin{equation}
\label{eq:pncon}
\begin{aligned}
&\mbox{(C1)}\quad N\circ\pi^\sharp=\pi^\sharp\circ N^\ast=:\pi^\sharp_N;%\nonumber
\\
&\mbox{(C2)}\quad[\alpha,\beta]_{\pi_N}=[N^\ast\alpha,\beta]_{\pi}+[\alpha,N^\ast\beta]_{\pi}-N^\ast[\alpha,\beta]_{\pi},\qquad\forall\alpha,\beta\in\Omega^1(M),
\end{aligned}
\end{equation}
where, if $P$ is a bivector field, 
the bracket $[-,-]_P:\Omega^1(M)\times\Omega^1(M)\rightarrow\Omega^1(M)$ is defined by 
\[
[\alpha,\beta]_P=\mathcal L_{P^\sharp\alpha}\beta
-\mathcal L_{P^\sharp\beta}\alpha-d\langle\beta,P^\sharp\alpha\rangle,
\]
and $\mathcal L_X$ is the Lie derivative with respect to the vector field $X$. A manifold $M$ endowed 
with a PN structure $(\pi,N)$ will be called a PN manifold and will be denoted by the triple $(M,\pi,N)$. 

\begin{remark}\label{rem:MMcon}
The second formula in 
\eqref{eq:pncon} is often written as $C(\pi,N)\equiv 0$, where 
\begin{equation}
C(\pi,N)(\alpha,\beta)=[\alpha,\beta]_{\pi_N}-[N^\ast\alpha,\beta]_{\pi}-[\alpha,N^\ast\beta]_{\pi}+N^\ast[\alpha,\beta]_{\pi},\qquad 
\alpha,\beta\in\Omega^1(M),\label{eq:MMconc}
\end{equation}
is the so called Magri-Morosi concomitant. It is also worth recalling that a torsion-free $(1,1)$ tensor field $N$ is often called a recursion operator since the vector space of all $X\in\mathfrak X(M)$ such that $\mathcal L_X N=0$ is an $N$-invariant Lie subalgebra of $\mathfrak X(M)$. This means, in particular, that if $X$ is an infinitesimal symmetry of $N$, i.e., 
$\mathcal L_XN=0$, then $N^kX$ will also be, for all $k\geq 1$.
\end{remark}
To every PN structure one can associate the space
\begin{equation}
\Omega^1_{Ham}(M)=\{\alpha\in\Omega^1(M)\mid d\alpha=0=d_N\alpha\},
\end{equation}
where $d_N:\Omega^\bullet(M)\rightarrow\Omega^{\bullet+1}(M)$ is defined by $d_N=i_N\circ d-d\circ i_N$ and $i_N$ is the degree-zero endormorphism of $\Omega^\bullet(M)$ defined on $\omega\in\Omega^k(M)$ by $i_N\omega(X_1,\dots,X_k)=\sum_{j=1}^k\omega(X_1,\dots,NX_j,\dots,X_k)$. The following proposition encloses an important property of $\Omega^1_{Ham}(M)$. 
\begin{prop}[\cite{Magri-Morosi,Bonechi}] The space $\Omega^1_{Ham}(M)$ is $N^\ast$-invariant. Moreover,
\begin{enumerate}
\item For every $k\geq 0$, it is closed with respect to $[-,-]_{\pi_{N^k}}$.
\item Given $\alpha\in\Omega^1_{Ham}(M)$ and $s\geq 1$, set $\alpha_s={N^\ast}^s\alpha$. Then 
\[
[\alpha_i,\alpha_j]_{\pi_{N^k}}=0
\] 
for all $k\geq 0$ and $i,j\geq 1$.
\end{enumerate}
\end{prop}
In other words, $\Omega^1_{Ham}(M)$ is a receptacle for the so called Lenard-Magri, or bi-Hamiltonian, chains.  
An important example is obtained as follows. For every $k\geq 1$, let $I_k:=\frac{1}{2k}\text{Tr}(N^k)$, where $\text{Tr}$ denote the trace operator. Then one can prove that $\{dI_k\}_{k\geq 1}\subset\Omega^1_{Ham}(M)$. More precisely, a direct computation shows that, for every $k\geq 1$, 
\begin{equation}
N^\ast dI_k=dI_{k+1},\label{eq:PNrec}
\end{equation}
which yields the fundamental involutivity property
\begin{equation}
\{I_i,I_j\}=0
\qquad\forall i,j\geq 1,\label{eq:invfam}
\end{equation}
where $\{-,-\}$ is the Poisson bracket defined by $\pi$ via the formula 
$\{f,g\}=\pi(df,dg)$, for all $f,g\in C^\infty(M)$.

\begin{example}[Das-Okubo PN structure \cite{DO}]
\label{ex:1}
On $M=\mathbb R^{2n}$, with coordinates $(p_i,q^i)$, $i=1,\dots,n$, let $\pi=\sum_{i=1}^{n}\frac{\partial}{\partial p_i}\wedge\frac{\partial}{\partial q^i}$ and 
\begin{equation}
\begin{split}
N_{DO}&=\sum_{i=1}^n p_i\Big(\frac{\partial}{\partial q^i}\otimes dq^i
+\frac{\partial}{\partial p_i}\otimes dp_i\Big)+\sum_{i<j}\Big(\frac{\partial}{\partial q^i}\otimes dp_j-\frac{\partial}{\partial q^j}\otimes dp_i\Big)\\
&+\sum_{i=1}^{n-1}e^{q^i-q^{i+1}}\Big(\frac{\partial}{\partial p_{i+1}}\otimes dq^i-\frac{\partial}{\partial p_i}\otimes dq^{i+1}\Big).\label{eq:DO-PN}
\end{split}
\end{equation}
Then $(\pi,N_{DO})$ is a PN structure and the sequence $\{I_k\}_{k\geq 1}$ is a set of first integrals for the \emph{open $n$-particle Toda chain}. In particular, $I_1$ is the total linear momentum and $I_2$ is the energy (Hamiltonian) of this system.
\end{example}

The notion of a Poisson quasi-Nijenhuis structure, introduced in \cite{SX}, represents an interesting generalization of the Poisson-Nijenhuis one. This enhanced setting is defined by a triple $(\pi,N,\phi)$, where $\pi$ is a Poisson tensor, 
$N$ is a tensor field of type (1,1) which is compatible with $\pi$, and $\phi$ is a $d$- and $d_N$-closed 3-form such that 
\begin{equation}
T_N(X,Y)=\pi^\sharp(\phi(X,Y,-))=:\pi^\sharp (i_{X\wedge Y}\phi)%,
\qquad\forall X,Y\in\mathfrak X(M).\label{eq:PqNtor}
\end{equation}
%\begin{remark}\label{rem:trivpqn}
%We will call the PqN structure $(\pi,N,\phi)$ \emph{trivial} if $\pi=0$.
%\end{remark}
Since $T_N$ does not vanish,  the corresponding space $\Omega_{Ham}^1(M)$ is not $N^\ast$-invariant and, for this reason, it looses the property of being the natural place where to look for Lenard-Magri chains. In particular, in this more general context, the recurrence equations \eqref{eq:PNrec} assume the form
\begin{equation}
dI_{k+1}=N^\ast dI_k+\phi_{k-1},\qquad\forall k\geq 1,\label{eq:genrec}
\end{equation}
where the 1-forms $\phi_s$ are defined by the formula
\begin{equation}
\langle\phi_s,X\rangle=\frac{1}{2}\text{Tr}(N^s(i_XT_N)),\qquad\forall X\in\mathfrak X(M), s\geq 0, \label{eq:phis}
\end{equation}
and $\langle i_XT_N,Y\rangle=T_N(X,Y)$ for all vector fields $Y$. The recurrence defined in \eqref{eq:genrec} entails that, for all $k>j\geq 1$,
\begin{equation}
\{I_k,I_j\}-\{I_{k-1},I_{j+1}\}=-\langle\phi_{j-1},\pi^\sharp dI_{k-1}\rangle-\langle\phi_{k-1},\pi^\sharp dI_j\rangle,\label{eq:recursionphi}
\end{equation}
which implies that, for a general PqN structure, the family $\{I_k\}_{k\geq 1}$ fails to be 
involutive, i.e., \eqref{eq:invfam} is not satisfied for all possible choices of $i,j$. This observation justifies the following

\begin{defi}\label{def:invpqn}
A PqN structure is called involutive if \eqref{eq:invfam} holds for all $i$ and $j$.
\end{defi}

Even though this definition is quite restrictive, the following example shows that there is at least one important class of PqN structures which meet its requirement.

\begin{example}[PqN structure of the closed Toda system]\label{ex:2}
An important example of involutive PqN structure 
was discovered in \cite{FMOP2020}, see also \cite{FMP2023}, and goes as follows. Let $(M,\pi)$ and $N_{DO}$ be as in Example \ref{ex:1}, and let
\begin{equation}
\label{PqN-Toda}
\widetilde{N}=N_{DO}-e^{q^n-q^1}\Big(\frac{\partial}{\partial p^1}\otimes dq_n-\frac{\partial}{\partial p_n}\otimes dq^1\Big).
\end{equation}
If $\phi=2e^{q^n-q^1}dq^1\wedge dq^n\wedge\sum_{i=1}^ndp_i$, then one can show that $(\pi,\widetilde N,\phi)$ is an involutive PqN structure on $M$. This family of PqN structures describes the so called \emph{closed $n$-particle Toda chain}. More precisely, and in analogy with the case of the Das-Okubo PN structure, the functions ${\tilde I}_1$ and 
${\tilde I}_2$ are, respectively, the total linear momentum and the energy (Hamiltonian) of this classical integrable system. 
\end{example}

It is therefore an interesting problem to pin down a set of assumptions entailing the involutivity of a PqN structure. 
A set of involutivity %sufficient 
conditions %ensuring the involutivity of a PqN structure 
was presented in \cite{FMOP2020} and subsequently improved in \cite{FMP2024}. We recall this result in 
\begin{theorem}\label{theo:inv} Let $(\pi,N,\phi)$ be a PqN structure on $M$ and let $I_k=\frac{1}{2k}\text{Tr}\left(N^k\right)$, $k\geq 1$. Suppose there is a 2-form $\Omega$ such that:
\begin{enumerate}
\item[(a)] $\phi=-2\,dI_1\wedge\Omega$;
\item[(b)] $\Omega(X_j,Y_k)=0$ for all $j,k\geq 1$, where $Y_k=N^{k-1}X_1-X_k$ and $X_k=\pi^\sharp dI_k$. 
\end{enumerate}
Then $\{I_j,I_k\}=0$ for all $j,k\ge 1$, where $\{-,-\}$ is the Poisson bracket defined by $\pi$.
\end{theorem}
Another involutivity theorem for PqN manifolds can be found in \cite{FMP2026}.

Before moving to the next section, we make a final comment about the family of PqN structures presented in Example \ref{ex:2}. To this end, we 
recall that 
a problem which is parallel, though independent, to the task of taming the class of involutive PqN structure, 
is the one of ``deforming" a PqN structure into another one. To phrase it in a more precise way, we recall the following theorem  \cite{FMOP2020,DMP2024}.

\begin{theorem}\label{thm:defthe} Let $(\pi,N,\phi)$ be a PqN structure on $M$ and let $\Omega$ be a closed 2-form. If 
\[
\widetilde N=N+\pi^\sharp\Omega^\flat\quad\text{and}\quad\widetilde\phi=\phi+d_N\Omega+
\frac{1}{2}[\Omega,\Omega]_{\pi},
\]
then $(\pi,\widetilde N,\widetilde\phi)$ is a PqN structure. In particular, if $\widetilde\phi=0$, 
the pair $(\pi,\widetilde N)$ is a PN structure.
\end{theorem} 
This result, together with Theorem \ref{theo:inv}, opens the door to the following problem: given a PqN structure, 
which assumptions on $\Omega$ guarantee that the deformed structure $(\pi,\widetilde N,\widetilde\phi)$ is involutive? Or, equivalently, when does a PqN structure admit involutive deformations? Note that, going back to the PN and PqN structures presented in Examples \ref{ex:1} and \ref{ex:2}, one can show that the latter can be obtained as a deformation of the former using $\Omega=e^{q^n-q^1}dq^n\wedge dq^1$, see \cite{FMOP2020,FMP2023,FMP2024}.

\section{Three-dimensional oriented PqN manifolds}\label{sec:3} 

In this section we will work with oriented manifolds, i.e., with pairs $(M,\mathcal V)$ where $M$ is a manifold and $\mathcal V$ is a never vanishing differential form of maximal degree, i.e., a volume form on $M$. Any such form defines a family of isomorphisms $\star_k:\Gamma(\Lambda^k(TM))\rightarrow\Omega^{\dim M-k}(M)$ via 
\[
\star_kP:=i_P\mathcal V,
\] 
that is, the contraction of $\mathcal V$ with $P$, which is the unique $(\dim M-k)$-form such that, for all 
$Q\in\Gamma(\Lambda^{\dim M-k}(M))$,
\[
\langle i_P\mathcal V,Q\rangle=\langle\mathcal V,P\wedge Q\rangle.
\] 
Furthermore, one can associate to every $X\in\mathfrak X(M)$  its divergence, i.e., the unique smooth function $\operatorname{div}(X)$ such that $\mathcal L_X\mathcal V=\operatorname{div}(X)\mathcal V$. Once this has been settled, we can start our study of  
three-dimensional oriented PqN manifolds, stating the following very well known result.

\begin{prop}[\cite{L-GPV}, Proposition 9.20]
\label{prop:xi1} 
Let $(M,\mathcal V)$ be a three-dimensional oriented manifold. 
A bivector field $\pi\in\Gamma(\Lambda^2TM)$ is Poisson if and only if $\xi\wedge d\xi=0$, 
where $\xi=\star_2\pi=i_\pi\mathcal V$.
In this case, 
\begin{equation}
\label{prop:xi2}
\pi^\sharp\xi=0.
\end{equation}
\end{prop}

\begin{remark}
\label{rem:vanish-xi-pi}
Since $\star_2$ is an isomorphism, the Poisson tensor $\pi$ never vanishes if and only if the corresponding 1-form $\xi$ never vanishes.
\end{remark}

We also recall a very nice characterization of three-dimensional oriented PN manifolds, which was presented in \cite{Cheng-Sheng}, Theorem 2.1. 

%\begin{remark} Hereafter we %will 
%limit ourself to consider the case of nowhere zero Poisson bivector fields, referring the reader to the above cited reference for an in-depth analysis of the general case.
%\end{remark}

\begin{theorem}\label{theo:charN}
Let $N$ be a $(1,1)$ tensor field on a three-dimensional oriented Poisson manifold $(M,\pi,\mathcal V)$. Suppose that $\pi$ never vanishes.
Then $(M,\pi,N)$ is a PN manifold if and only if $N$ can be decomposed as 
\begin{equation}
N=\lambda I+Z\otimes\xi,\label{eq:chPN}
\end{equation}
where $\lambda\in C^\infty(M)$, $Z\in\mathfrak X(M)$, $\xi=i_\pi\mathcal V$ and the following conditions hold:
\begin{equation}
d\left(\lambda+\left<\xi,Z\right>\right)=\operatorname{div}(Z)\xi\quad\text{and}\quad Z(\lambda)=0.
\label{eq:contor}
\end{equation}
%\begin{equation}
%\xi\wedge(d\left(\left<\xi,Z\right>+\lambda\right))=0\quad\text{and}\quad Z(\lambda)=0.\label{eq:contor}
%\end{equation}
\end{theorem}

A couple of remarks about the proof of the previous statement are in order.
\begin{remark}\label{rem:2stat}
\begin{enumerate}
\item[]
\item Under the hypotheses of the previous theorem, the condition $N\circ\pi^\sharp=\pi^\sharp\circ N^\ast$ forces the decomposition of the $(1,1)$ tensor field $N$ as $N=\lambda I+Z\otimes\xi$, where $\lambda$, $Z$ and $\xi$ are as in the statement of the theorem, see Lemma 2.4 in \cite{Cheng-Sheng}. On the other hand, 
since $\pi^\sharp\xi=0$ by (\ref{prop:xi2}), if $N=\lambda I+Z\otimes\xi$ then $N\circ\pi^\sharp=\pi^\sharp\circ N^\ast$. For this reason one has that
\[
N=\lambda I+Z\otimes\xi\iff N\circ\pi^\sharp=\pi^\sharp\circ N^\ast.
\]

Furthermore, given the above decomposition of $N$, 
\[
C(\pi,N)\equiv 0 \iff d(\lambda+\langle\xi,Z\rangle)=\operatorname{div}(Z)\xi,
\] 
see Lemma 2.6 in \cite{Cheng-Sheng}. In other words, on an oriented three-dimensional manifold, the pair $(\pi,N)$ is compatible, see \eqref{eq:pncon}, if and only if
\[
N=\lambda I+Z\otimes\xi\quad\text{and}\quad d(\lambda+\langle\xi,Z\rangle)=\operatorname{div}(Z)\xi.
\]
\item On the other hand, in Lemma 2.5 of \cite{Cheng-Sheng} it is proven that if $N$ decomposes as above, then 
\begin{equation}
T_N\equiv 0\iff 
\xi\wedge d(\lambda+\langle\xi,Z\rangle)= 0\quad\text{and}\quad Z(\lambda) = 0.\label{eq:equiv}
\end{equation}
If $\pi$ and $N$ are compatible, then the first identity in the right-hand side of the previous equivalence is automatically verified, so one obtains the proof of Theorem \ref{theo:charN}. Note that  
\eqref{eq:equiv} can be proved using the identity
\begin{equation}\label{eq: T_N computation} %%%%%%%%%%%
T_N(X,Y)=\left(\xi\wedge d\left(\lambda+\left<\xi,Z\right>\right)\right)(X,Y)Z+Z(\lambda)i_{\xi}\left(X\wedge Y\right),
\quad\forall X,Y\in\mathfrak X(M),
\end{equation}
see the proof of Lemma 2.5 in \cite{Cheng-Sheng}.
\end{enumerate}
\end{remark}

We are now ready to state the main result of this section which provides %give 
necessary and sufficient conditions for a triple $(\pi,N,\phi)$, with $\pi$ never vanishing, to be a PqN structure on an oriented three-dimensional manifold $(M,\mathcal V)$.

%\begin{prop}
\begin{theorem}\label{prop 3D PqN} %%%%%%%%%%%%
Let $(M,\pi)$ be an oriented three-dimensional manifold with volume form $\mathcal V$. 
Suppose that $\pi$ never vanishes.
Then the quadruple $(M,\pi,N,\phi)$ is a PqN manifold if and only if $N$ can be decomposed as $N=\lambda I+Z\otimes\xi$, where $\xi=i_{\pi}\mathcal V$,  $Z\in\mathfrak X(M)$, $\lambda\in C^{\infty}(M)$ and the following conditions hold:
\begin{equation}\label{eq: 3D PqN conditions}
d\left(\lambda+\left<\xi,Z\right>\right)=\operatorname{div}(Z)\xi\quad\text{and}\quad\phi=-Z(\lambda)\mathcal V.
\end{equation}
\end{theorem}
%\end{prop}

\begin{proof} As seen in the first item of Remark \ref{rem:2stat}, 
the pair $(\pi,N)$ is compatible if and only if
\[
N=\lambda I+Z\otimes\xi\quad\text{and}\quad d(\lambda+\langle\xi,Z\rangle)=\operatorname{div}(Z)\xi.
\]
It remains to prove that $\phi=-Z(\lambda)\mathcal V$. To this end, we first observe that \eqref{eq: T_N computation} yields 
\begin{equation*}
T_N(X,Y)=\cancel{\left(\xi\wedge \operatorname{div}(Z)\xi\right)}(X,Y)Z+Z(\lambda)i_{\xi}(X\wedge Y)
\end{equation*}
i.e.,
\begin{equation}
T_N(X,Y)=Z(\lambda)i_\xi(X\wedge Y).\label{eq: T_N 3D PqN}
\end{equation}
Now, for every $\eta\in\Omega^1(M)$, 
\begin{align*}
\left<\eta,\pi^{\sharp}(i_{X\wedge Y}\phi)\right>&=\pi\left(i_{X\wedge Y}\phi,\eta\right)=\left<\left(i_{X\wedge Y}\phi\right)\wedge \eta,\pi\right>
=\left<\left(i_Y\left(i_X\phi\right)\right)\wedge\eta,\pi\right>\\
&=
\left<i_Y\left(\left(i_X\phi\right)\wedge\eta\right)-\left(i_X\phi\right)\wedge \left(i_Y\eta\right),\pi\right>
=\left<i_Y\left(\left(i_X\phi\right)\wedge\eta\right),\pi\right>-\left<\eta,Y\right>\left<i_X\phi,\pi\right>
\\
&=\left<i_Y\left(i_X\left(\cancel{\phi\wedge\eta}\right)+\phi\wedge\left(i_X\eta\right)\right),\pi\right>-\left<\eta,Y\right>\left<i_X\phi,\pi\right>\\
&=\left<\eta,X\right>\left<i_Y\phi,\pi\right>-\left<\eta,Y\right>\left<i_{\pi}\phi,X\right>=\left<\eta,X\right>\left<i_{\pi}\phi,Y\right>-\left<\eta,Y\right>\left<i_{\pi}\phi,X\right>\\
&=\left<\eta,\left<i_{\pi}\phi,Y\right>X-\left<i_{\pi}\phi,X\right>Y\right>
=-\left<\eta,i_{i_{\pi}\phi}(X\wedge Y)\right>,
\end{align*}
so that 
\begin{equation}
\label{eq:iiphi}
\pi^{\sharp}(i_{X\wedge Y}\phi)=-i_{i_{\pi}\phi}(X\wedge Y).
\end{equation}
Thanks to \eqref{eq: T_N 3D PqN} and \eqref{eq:iiphi}, we have that $T(N)(X,Y)=\pi^{\sharp}(i_{X\wedge Y}\phi)$ 
if and only if $i_{i_{\pi}\phi}(X\wedge Y)=-i_{Z(\lambda)\xi}(X\wedge Y)$. Since $X,Y\in\mathfrak X(M)$ are arbitrary, this is equivalent to 
$$i_{\pi}\phi=-Z(\lambda)\xi.$$ But $\phi=f\mathcal V$ for some $f\in C^\infty(M)$, and 
$\xi=i_\pi\mathcal V$, so we arrive at the condition $(f+Z(\lambda))\xi=0$, which is equivalent to $f=-Z(\lambda)$ since $\xi$ never vanishes --- see Remark \ref{rem:vanish-xi-pi}.
We can therefore conclude that $T_N(X,Y)=\pi^{\sharp}(i_{X\wedge Y}\phi)$ if and only if $\phi=-Z(\lambda)\mathcal V$.
\end{proof}

\begin{remark}\label{rem:indpqn} %%%%%%%%%%%
It is worth making a comment about the independence of the PqN structure from the chosen volume form $\mathcal V$. First note that if $\mathcal V'$ is another volume form, then $\mathcal V'=f\mathcal V$ for a given never-vanishing smooth function $f$. This implies the equalities $\xi'=f\xi$, $\lambda'=\lambda$, and $Z'=\frac{1}{f}Z$, so that $d(\lambda+\left<\xi,Z\right>)=d(\lambda'+\left<\xi',Z'\right>)$ and 
${\operatorname{div}}_{\mathcal V}\left(Z\right)\xi={\operatorname{div}}_{\mathcal V'}\left(Z\right)\xi'$, see Remark 2.3 in \cite{Cheng-Sheng}. Finally, 
\[
\phi'=-Z'(\lambda)\mathcal V'=-Z(\lambda)\mathcal V=\phi,
\] 
proving that the PqN-structure is independent of the choice of the volume form.
\end{remark} 

\begin{example}\label{exa:PqN-R3} %%%%%%%%%%%
We present a family of PqN structures on $\mathbb R^3$, to be compared with the corresponding family of PN 
structures in \cite{Cheng-Sheng}, Proposition 3.1. Let $\pi=\frac{\partial}{\partial x}\wedge\frac{\partial}{\partial y}\in\Gamma(\Lambda^2 T\mathbb R^3)$ 
and let $\mathcal V=dx\wedge dy\wedge dz$ be the canonical volume form on $\mathbb R^3$. 
If $N$ is a (1,1) tensor field and $\phi$ is a 3-form, we know from %Proposition 
Theorem \ref{prop 3D PqN}
that the quadruple $(\mathbb R^3,\pi,N,\phi)$ is a PqN manifold if and only if 
$N=\lambda I+Z\otimes\xi$, where $\xi=i_{\pi}\mathcal V=dz$, and the following conditions hold:
\begin{equation*}
d\left(\lambda+\left<\xi,Z\right>\right)=\operatorname{div}(Z)\xi\quad\text{and}\quad\phi=-Z(\lambda)\mathcal V.
\end{equation*}
If 
\begin{equation}
\label{eq: Z}
Z=a\frac{\partial}{\partial x}+b\frac{\partial}{\partial y}+c\frac{\partial}{\partial z},
\end{equation}
then 
\begin{align}
d(\left<\xi,Z\right>+\lambda)=
\operatorname{div}(Z)
\xi\ \Longleftrightarrow\ (c+\lambda)_xdx+(c+\lambda)_ydy+(c+\lambda)_zdz
=(a_x+b_y+c_z)dz,\label{eq:equico}
\end{align}
which is equivalent to
\begin{align*}
\begin{cases}
(c+\lambda)_x=0\\
(c+\lambda)_y=0\\
\lambda_z=a_x+b_y.
\end{cases}
\end{align*}
Thus $c+\lambda$ is a function of $z$ only, say, $g(z)$. 
Since $Z(\lambda)=a\lambda_x+b\lambda_y+c\lambda_z$, one has that
\begin{equation}
\phi=-Z(\lambda)\mathcal V=(ac_x+bc_y-c(a_x+b_y))dx\wedge dy\wedge dz.\label{eq:phi}
\end{equation}
Hence we found a recipe to complete $\pi=\frac{\partial}{\partial x}\wedge\frac{\partial}{\partial y}$ 
to a PqN structure on $\mathbb R^3$ endowed with the volume form $\mathcal V=dx\wedge dy\wedge dz$. 
One has to choose arbitrary functions $\lambda,%\,
a\in C^{\infty}(\mathbb{R}^3)$ and $g(z)$. Then, after defining
\[
b=\int({\lambda}_z-a_x)dy\quad\text{and}\quad c=g-\lambda,
\]
one constructs the vector field (\ref{eq: Z}) and the PqN structure $(\mathbb{R}^3,\pi,N,\phi)$, where $N=\lambda I+Z\otimes dz$ and $\phi$ is given by \eqref{eq:phi}.
\end{example}

%Let  $(\pi,N,\phi)$ be a PqN structure. To notice whether such PqN structure is a deformation obtained from  PN structure, in the sense of Theorem \ref{thm:defthe}, we must %determine if there exists
%verify the existence of  a closed 2-form $\Omega$ %verifying:
%such that
%\begin{itemize}
%\item the (1,1) tensor field 
%\begin{equation}\label{eq: N_1}
%N_1:=N-\pi^{\sharp}\Omega^{\flat}
%\end{equation} 
%is a Nijenhuis operator compatible with $\pi$;
%\item  $\Omega$ and $\phi$ are related by
%\begin{equation}
%\phi=d_{N_1}\Omega+\frac{1}{2}\left[\Omega,\Omega\right]_{\pi}.\label{eq:phi}
%\end{equation}
%\end{itemize} 
%Under these conditions the pair $(\pi,N_1)$ is a PN manifold and its deformation by $\Omega$ is the PqN structure $(\pi,N,\phi)$. 
%The following proposition shows that  
% this can be always done for an 
A first application of Theorem \ref{prop 3D PqN} focuses on the problem of deformed PqN structures in the sense of Theorem \ref{thm:defthe}. 
%every oriented three-dimensional PqN structure with a never vanishing Poisson bivector can be obtained, locally, as a deformation of a PN manifold, see Theorem \ref{thm:defthe}.
\begin{prop}\label{prop:local}
Locally, every PqN structure on an oriented three-dimensional manifold with a never vanishing Poisson bivector is obtained by deforming a PN structure by means of a closed 2-form. 
\end{prop}
\begin{proof}
Let $(\pi,N,\phi)$ be a PqN structure, with $\pi$ never vanishing, on an oriented three-dimensional manifold  $\left(M,\mathcal V\right)$ and let $m\in M$. Darboux theorem entails that there exists a coordinate neighborhood $U$ of  $m$ with coordinates, say $(x,y,z)$, such that $\pi=\frac{\partial}{\partial x}\wedge\frac{\partial}{\partial y}$ on $U$. %The independence of the PqN structure with respect to the volume form, see Remark \ref{rem:indpqn}, allows us to consider $\mathcal V=dx\wedge dy\wedge dz$ on $U$. 
Moreover, one can consider $\mathcal V=dx\wedge dy\wedge dz$ on $U$, see Remark \ref{rem:indpqn}. Then, on $U$, Example \ref{exa:PqN-R3} shows that $N=\lambda I+Z\otimes \xi$ where $\xi=dz$ and $Z$ is given by \eqref{eq: Z} verifying $\lambda+c=g(z)$ and $\lambda_z=a_x+b_y$.  Now consider
\[
\Omega=\Omega_{12}dx\wedge dy+\Omega_{13}dx\wedge dz+\Omega_{23}dy\wedge dz
=-\lambda\, dx\wedge dy-b\,dx\wedge dz+a\,dy\wedge dz,
\]
and note that $d\Omega=0$ since $\lambda_z=a_x+b_y$. Using $\lambda+c=g(z)$, %we compute the $(1,1)$ tensor  field $N_1:=N-\pi^{\sharp}\Omega^{\flat}$: 
one can compute
\begin{align*}
\pi^{\sharp}\Omega^{\flat}%&=-\Omega_{12}\frac{\partial}{\partial x}\otimes dx-\Omega_{12}\frac{\partial}{\partial y}\otimes dy+\Omega_{23}\frac{\partial}{\partial x}\otimes dz-\Omega_{13}\frac{\partial}{\partial y}\otimes dz\\
&=-\Omega_{12} I+\left(\Omega_{23}\frac{\partial}{\partial x}-\Omega_{13}\frac{\partial}{\partial y}+\Omega_{12}\frac{\partial}{\partial z}\right)\otimes dz\\
&=\lambda I+\left(a\frac{\partial}{\partial x}+b\frac{\partial}{\partial y}-(g(z)-c)\frac{\partial}{\partial z}\right)\otimes dz\\
&=\lambda I+%(a\partial_x+b\partial_y+c\partial_z)
Z\otimes dz-g(z)\frac{\partial}{\partial z}\otimes dz.
\end{align*}
%yielding the identity
%\begin{equation}
%g(z)\frac{\partial}{\partial z}\otimes dz+\pi^{\sharp}\Omega^{\flat}=N.\label{eq:identN}
%\end{equation}
If $\widetilde{Z}:=g(z)\frac{\partial}{\partial z}$ and $N_1:=\widetilde{Z}\otimes dz$, 
then %one can write \eqref{eq:identN} in the form 
$N=N_1+\pi^{\sharp}\Omega^{\flat}$.
%\begin{align*}
%N_1&=(\lambda I+Z\otimes dz)-\left(-\Omega_{12}\frac{\partial}{\partial x}\otimes dx-\Omega_{12}\frac{\partial}{\partial y}\otimes dy+\Omega_{23}\frac{\partial}{\partial x}\otimes dz-\Omega_{13}\frac{\partial}{\partial y}\otimes dz\right)\\
%&=(\lambda I+Z\otimes dz)+\Omega_{12} I-\left(\Omega_{23}\frac{\partial}{\partial x}-\Omega_{13}\frac{\partial}{\partial y}+\Omega_{12}\frac{\partial}{\partial z}\right)\otimes dz\\
%&=Z\otimes dz-\left(a\frac{\partial}{\partial x}+b\frac{\partial}{\partial y}-(g(z)-c)\frac{\partial}{\partial z}\right)\otimes dz=g(z)\frac{\partial}{\partial z}\otimes dz.
%\end{align*}
%If we let $\widetilde{Z}:=g(z)\frac{\partial}{\partial z}$ and $N_1:=\widetilde{Z}\otimes dz$ then $N=N_1+\pi^{\sharp}\Omega^{\flat}$. 
One can observe that $N_1$ is a $(1,1)$ tensor field compatible with $\pi=\frac{\partial}{\partial x}\wedge\frac{\partial}{\partial y}$, since $N_1\circ\pi^{\sharp}=\pi^{\sharp}\circ N_1^{*}$ and
\[
d\langle \xi,\widetilde{Z}\rangle=d(g(z))=g'(z)dz=\operatorname{div}(\widetilde{Z})\xi,
\]
see Remark \ref{rem:2stat}. Moreover, \eqref{eq:equiv} shows that $T_{N_1}=0$, proving that $(\pi,N_1)$ is a PN structure. 
To conclude that the original PqN structure is a deformation of this PN structure, we are left to show 
\begin{equation}
d_{N_1}\Omega+\frac{1}{2}\left[\Omega,\Omega\right]_{\pi}=\phi=-Z(\lambda)\mathcal V,\label{eq:dNphi}
\end{equation}
whose computational details are enclosed in Section \ref{subsec:app1}.
\end{proof}

Another significant application of Theorem \ref{prop 3D PqN}  is its use in the following
%\begin{theorem}
\begin{prop}\label{thm:inv3}
Every three-dimensional oriented PqN manifold with a never vanishing Poisson tensor is involutive.
\end{prop}
\begin{proof}
Using local Darboux coordinates $(x,y,z)$ and the notation as in Example \ref{exa:PqN-R3}, one has
\begin{equation}
N=\begin{bmatrix}\lambda(x,y,z) & 0 & a(x,y,z)\\ 0 & \lambda(x,y,z) & b(x,y,z) \\ 0 & 0 & g(z)\end{bmatrix}.\label{eq:thanksRef}
\end{equation}
This implies that $\text{Tr} N^k=2\lambda^k+g^k$ for every $k\geq 1$. Since $\pi=\frac{\partial}{\partial x}\wedge\frac{\partial}{\partial y}$, 
one has that 
\[
\{I_i,I_j\}=\pi(dI_i,dI_j)=0
\]
for all $i,j$.\end{proof}

For an alternative proof of this proposition which does not use local coordinates, we refer the reader to Subsection \ref{ss:proofth}.

%The proof of this statement is based on a closed formula for the 1-forms $\phi_k$, see \eqref{eq:genrec}, which is contained in the following

%\begin{lem}\label{lem:phik}
%On an oriented three-dimensional PqN manifold, where $N=\lambda I+Z\otimes\xi$, we have that 
%\begin{equation}
%\phi_k=Z(\lambda)\lambda^k\xi\qquad\mbox{for all $k\geq 0$.}
%\label{eq:phik}
%\end{equation}
%\end{lem}
%\begin{proof}
%By induction on $k$, one proves (see Section \ref{subsec:app2}) that 
%\begin{equation}
%N^k=\lambda^kI+f_k \,Z\otimes\xi.\label{eq:Nalk}
%\end{equation}
%where $f_k=\sum_{l=0}^{k-1}{k\choose l}\lambda^l\langle\xi,Z\rangle^{k-l-1}$. Moreover, since
%\[
%T_N(X,Y)\stackrel{\eqref{eq: T_N 3D PqN}}{=}Z(\lambda)i_\xi(X\wedge Y)=Z(\lambda)\langle\xi,X\rangle Y-Z(\lambda)\langle\xi,Y\rangle X,
%\]
%one has that
%\begin{equation}
%i_XT_N=Z(\lambda)(\langle\xi,X\rangle I-X\otimes\xi).\label{eq:iXTN}
%\end{equation}
%For every $k\geq 0$ and $X\in\mathfrak X(M)$, we can now compute
%\[
%\begin{split}
%\langle\phi_k,X\rangle&\stackrel{%\eqref{eq:phi}
%\eqref{eq:phis}}{=}\frac12\text{Tr}(N^ki_XT_N)\stackrel{\eqref{eq:iXTN}}
%{=}\frac12\text{Tr}\big(N^k(Z(\lambda)\langle\xi,X\rangle I-Z(\lambda)X\otimes\xi)\big)\\
%&\stackrel{\eqref{eq:Nalk}}{=}\frac12\text{Tr}\Big[\Big(\lambda^kI+\sum_{l=0}^{k-1}{k\choose %l}\lambda^l\langle\xi,Z\rangle^{k-l-1}Z\otimes\xi \Big)\Big(Z(\lambda)\langle\xi,X\rangle %I-Z(\lambda)X\otimes\xi)\Big)\Big]\\
%&=\frac12\text{Tr}\Big[Z(\lambda)\lambda^k\langle\xi,X\rangle %I-Z(\lambda)\lambda^kX\otimes\xi+Z(\lambda)\langle\xi,X\rangle\sum_{l=0}^{k-1}{k\choose %l}\lambda^l\langle\xi,Z\rangle^{k-l-1}Z\otimes\xi\\
%&-Z(\lambda)\sum_{l=1}^{k-1}{k\choose l}\lambda^l\langle\xi,Z\rangle^{k-l-1}\langle\xi,X\rangle Z\otimes\xi\Big],
%\end{split}
%\]
%where, in the last sum, $\langle\xi,X\rangle Z\otimes\xi$ comes from the  computation
%\[
%\left((Z\otimes\xi)\circ %(X\otimes\xi)\right)(Y)=\langle\xi,Y\rangle(Z\otimes\xi)(X)=\langle\xi,Y\rangle\langle\xi,X\rangle Z%,
%\qquad\forall Y\in\mathfrak X(M),
%\]
%which yields the identity $(Z\otimes\xi)\circ(X\otimes\xi)=\langle\xi,X\rangle Z\otimes\xi$. Now, recalling that $\text{Tr}(Z\otimes\xi)=\langle\xi,Z\rangle$ and $\text{Tr}(I)=3$, the last term of the previous chain of identities becomes
%\[
%\begin{split}
%&\frac12\left(3Z(\lambda)\lambda^k\langle\xi,X\rangle-\lambda^kZ(\lambda)\langle\xi,X\rangle+\cancel{Z(\lambda)\langle\xi%,X\rangle\sum_{l=0}^{k-1}{k\choose %l}\lambda^l\langle\xi,Z\rangle^{k-l}}\cancel{-Z(\lambda)\langle\xi,X\rangle\sum_{l=0}^{k-1}{k\choose %l}\lambda^l\langle\xi,Z\rangle^{k-l}}
%\right)\\
%&=Z(\lambda)\lambda^k\langle\xi,X\rangle,
%\end{split}
%\]
%which entails $\phi_k=Z(\lambda)\lambda^k\xi$.
%\end{proof}
%With this result at hand, we can go back to the
%\begin{proof}(of %Theorem
%Proposition \ref{thm:inv3}) We want to prove that $\{I_k,I_j\}=0$ for all $k,j$ with $k>j$. 
%Since the 1-forms $\phi_k$ are multiple of $\xi$, see \eqref{eq:phik}, and $\pi^\sharp\xi=0$, see (\ref{prop:xi2}),
%the right-hand side of \eqref{eq:recursionphi} is %(identically) 
%zero. In other words, that identity collapses to
%\[
%\{I_k,I_j\}=\{I_{k-1},I_{j+1}\}.
%\] 
%If $k-1=j+1$ we are done, otherwise one can apply again this argument, which, after a finite number of steps, will %entail the thesis.
%\end{proof}

%\begin{remark} 
%\label{rem:torsion-of-k-power}
%Using (\ref{eq: T_N computation}) and (\ref{eq:Nalk}), one obtains the following 
%generalization of (\ref{eq: T_N 3D PqN}), giving the Nijenhuis torsion of the $k$-th power of the tensor 
%$N=\lambda I+Z\otimes\xi$:
%\begin{equation}
%T_{N^k}(X,Y)=f_k Z(\lambda^k)i_\xi(X\wedge Y).\label{eq: T_N 3D PqN-k}
%\end{equation}
%This could also be deduced from the general expression in Remark 3 of \cite{BogoTens}.
%\end{remark}

%We close this section with the following remark, where we compare the result just obtained with Theorem \ref{theo:inv}, giving sufficient conditions for the involutivity of a (general) PqN manifold.
%\begin{remark} \label{rem:FMP2024}
%Applying $\pi^\sharp$ to both sides of \eqref{eq:genrec} and using again (\ref{prop:xi2}) and Lemma \ref{lem:phik}, one arrives to the identity $X_{k+1}=N X_k$, which holds for all $k\geq 1$ and which, if iterated, gives $X_{k+1}=N^{k}X_1$ for all $k\geq 1$.
%This observation implies that the vector fields $Y_k$ appearing in item (b) of Theorem \ref{theo:inv} are identically zero. 
%On the other hand, one can find a 2-form $\Omega$ such that $\phi=-Z(\lambda)\mathcal V=-2dI_1\wedge\Omega$, see for example \cite{Pham}, Chapter 3, Section 1.1. From these observations we can deduce that both conditions of Theorem \ref{theo:inv} are satisfied, and we obtain another proof of %Theorem
%Proposition \ref{thm:inv3}.  
%\end{remark}


\section{PqN and Haantjes structures}\label{sec:4}

We start this section recalling the definition of Haantjes torsion of a tensor field of type $(1,1)$, then we introduce the notion of Haantjes manifold and we discuss how this relates to the geometry of PqN manifolds. 

Let $N$ be a (1,1) tensor field. Recall that its Nijenhuis torsion $T_N$ is the tensor field of type $(1,2)$, i.e., the 2-form with values in $TM$, defined in \eqref{eq:torsionA}. To address the problem of the Frobenius integrability of 2-dimensional planes generated by the eigenvectors of $N$ (see Subsection \ref{ss:haantjes} 
and the references therein), 
Haantjes introduced in \cite{Haantjes} what is now known as the \emph{Haantjes torsion} of $N$. This is the tensor field 
$H_N$ of type $(1,2)$ defined by the formula
\begin{equation}
H_N(X,Y)=T_N(NX,NY)-N(T_N(NX,Y)+T_N(X,NY)-NT_N(X,Y)),\qquad
X,Y\in\mathfrak X(M), \label{eq:haant}
\end{equation}
where $T_N$ is the Nijenhuis torsion of $N$. The tensor $N$ will be called a Haantjes operator if $H_N\equiv 0$.

In \cite{BogoTens}, Bogoyavlenskij derived algebraic identities for Nijenhuis and Haantjes torsions on an arbitrary manifold. some of these results, which will be used subsequently, are collected in the following
\begin{lem}[\cite{BogoTens}, Corollaries 3, 4]\label{lem:Bogo}
Let  $N$ be a (1,1) tensor field on a %$n$-dimensional 
manifold $M$ and let $p(x,t)=\sum_{m=0}^ka_m(x)t^m$ be a polynomial with coefficients in $C^\infty(M)$.  Then: %We have three statements
\begin{enumerate}
\item If $N$ is a Haantjes operator, $p(x,N)$ is also a Haantjes operator.
\item For all $f,\,g\in C^{\infty}(M)$
\begin{equation}
H_{fI+gN}(X,Y)=g^4H_{N}(X,Y).\label{eq: Haantjes_defor_I}
\end{equation}
%where $I$ is the identity operator on $\mathfrak{X}(M)$.
\end{enumerate}
\end{lem}
The next example presents a formula for the Haantjes torsion of %a (1,1) tensor field which is 
the tensor product of a vector field and a 1-form.
\begin{example} 
Let $N=W\otimes\eta$, where $W\in\mathfrak{X}(M)$ and $\eta\in\Omega^{1}(M)$. Then for all $X,Y\in\mathfrak{X}(M)$
\begin{equation}
T_{N}(X,Y)=\left[\left<\eta,X\right>Y(\left<\eta,W\right>)-\left<\eta,Y\right>X(\left<\eta,W\right>)-(\eta\wedge d\eta)(W,X,Y)\right]W,\label{eq: NtorsionWxn}
\end{equation}
see Lemma 2.5 in \cite{Cheng-Sheng}.
Moreover a computation using \eqref{eq: NtorsionWxn} shows that
\begin{align*}
T_N(NX,NY)&=0,\\
NT_N(NX,Y)&=\left<\eta,W\right>\left[\left<\eta,X\right>\left<\eta,W\right>Y(\left<\eta,W\right>)-\left<\eta,Y\right>\left<\eta,X\right>W(\left<\eta,W\right>)\right]W,\\
NT_N(X,NY)&=\left<\eta,W\right>\left[\left<\eta,X\right>\left<\eta,Y\right>W(\left<\eta,W\right>)-\left<\eta,Y\right>\left<\eta,W\right>X(\left<\eta,W\right>)\right]W,\\
N^2T_N(X,Y)&=\left<\eta,W\right>^2\left[(\left<\eta,X\right>Y(\left<\eta,W\right>)-\left<\eta,Y\right>X(\left<\eta,W\right>)-(\eta\wedge d\eta)(W,X,Y)\right]W.
\end{align*}
Summing up the above four equations, \eqref{eq:haant} yields
\begin{equation}\label{eq: HtorsionWxn}
H_{N}(X,Y)=-\left<\eta,W\right>^2(\eta\wedge d\eta)(W,X,Y)W.
\end{equation}
\end{example}

As we already observed, the presence of a PN or, more generally, of a bi-Hamiltonian structure on $M$, makes it possible to introduce a class of distinguished vector fields (called bi-Hamiltonian), which possess a large number of conserved quantities which are mutually in involution with respect to the underlying Poisson structures. In particular, in the PN case, the presence of the recursion operator defines a natural space where to look for the Lenard-Magri chains. In the attempt to weaken the hypothesis on the recursion operator to be torsion free, but still keeping the possibility to define a notion of Lenard-Magri chain, Magri introduced in \cite{MagriHaa} the notion of Haantjes manifold, which we recall in %here below

\begin{defi} A Haantjes structure on a manifold $M$ is a pair $(N,\theta)$ of a $(1,1)$ tensor field $N$ 
and a 1-form $\theta$ such that:
\begin{itemize}
\item[(H1)] $H_N=0$;
\item[(H2)] $d\theta=0$;
\item[(H3)] $d_N\theta=0$;
\item[(H4)] $\langle\theta,T_N(X,Y)\rangle=0$ for all $X,Y\in\mathfrak X(M)$.
\end{itemize}
A manifold endowed with a Haantjes structure is called a Haantjes manifold. If $\theta=0$, the corresponding Haantjes structures is called trivial.
\end{defi}

We can now state and prove the main result of this section. 
\begin{theorem}
\label{theo:haan} 
If $(\pi,N,\phi)$ is a PqN structure on an oriented three-dimensional manifold $(M,\mathcal V)$ and $\pi$ never vanishes, then  $M$ admits a Haantjes structure $(N,\theta)$. 
\end{theorem}
\begin{proof} 
We know from Theorem \ref{prop 3D PqN} that 
$N=\lambda I+Z\otimes\xi$, where $\xi=i_{\pi}\mathcal{V}$, $\lambda\in C^{\infty}(M)$ and $Z\in\mathfrak X(M)$ are related by the first equation of \eqref{eq: 3D PqN conditions}. We claim that the pair $(N,\theta)$, where 
\begin{equation}\label{eq: 1-form theta}
\theta=\text{div}(Z)\xi,
\end{equation}
is a Haantjes structure on $M$. 
The exactness of $\theta$, see \eqref{eq: 3D PqN conditions}, gives us (H2). Now observe that, for every $X\in\mathfrak X(M)$,
\[
i_N\theta(X)=\langle\theta,NX\rangle=\langle\theta,\lambda X+\langle\xi,X\rangle Z\rangle=\lambda\langle\theta,X\rangle+\langle\theta,Z\rangle\langle\xi,X\rangle=(\lambda+\langle\xi,Z\rangle)\langle\theta,X\rangle,
\]
which yields the identity
\begin{equation}
i_N\theta=(\lambda+\langle\xi,Z\rangle)\theta.\label{eq:iNxi}
\end{equation}
Since $d\theta=0$, one has that
\[
d_N\theta=-d(i_N\theta)\stackrel{\eqref{eq:iNxi}}{=}-d((\lambda+\langle\xi,Z\rangle)\theta)=-d(\lambda+\langle\xi,Z\rangle)\wedge\text{div}(Z)\xi\stackrel{\eqref{eq: 3D PqN conditions}}{=}-\left(\text{div}(Z)\right)^2(\xi\wedge\xi)=0,
\]
proving (H3). Moreover, 
\[
\langle\theta,T_N(X,Y)\rangle\stackrel{\eqref{eq: T_N 3D PqN}}{=}Z(\lambda)\text{div}(Z)\langle\xi,i_\xi(X\wedge Y)\rangle=Z(\lambda)\text{div}(Z)\langle\xi,\langle\xi,X\rangle Y-\langle\xi,Y\rangle X\rangle=0
\]
for all $X,Y\in\mathfrak X(M)$, proving (H4). We are left to prove (H1). 
Applying \eqref{eq: Haantjes_defor_I} to $N=\lambda I+Z\otimes\xi$ and using \eqref{eq: HtorsionWxn}, we obtain
\[
H_N(X,Y)=H_{Z\otimes\xi}(X,Y)=-\left<\xi,Z\right>^2\left(\xi\wedge d\xi\right)(Z,X,Y)Z,
\]
which vanishes identically, see Proposition \ref{prop:xi1}.
\end{proof}

\begin{remark}
\begin{enumerate}
\item[]
\item In the same hypotheses of the previous theorem, Remark \ref{rem:vanish-xi-pi} implies that 
the Haantjes structure $(N,\theta)$ is non-trivial if the divergence of the vector field $Z$ is not zero.
\item For a general PqN structure $(\pi,N,\phi)$, the Haantjes torsion of $N$ does not vanish. An example is given by the case of the periodic Toda lattice, see \eqref{PqN-Toda}.
\end{enumerate}
\end{remark}

\begin{remark}
\label{rem:Eber}
A Haantjes structure $(N,\theta)$ on an oriented three-dimensional manifold cannot always be derived from a PqN structure $(N,\pi,\phi)$, since a non-trivial Poisson bivector $\pi$ such that  $N\circ\pi^{\sharp}=\pi^{\sharp}\circ N^*$ does not always exist. Indeed, consider 
(see Example 2 in \cite{MagriVeselov}) the (1,1) tensor field
\[
N=-\frac{y}{2}\frac{\partial}{\partial x}\otimes dz+2\frac{\partial}{\partial y}\otimes dx-z\frac{\partial}{\partial y}\otimes dz+2\frac{\partial}{\partial z}\otimes dy.
\]
It can be checked that its Haantjes torsion vanishes and that its Nijenhuis torsion
\[
T_N=\frac{\partial}{\partial x}\otimes dz\otimes dx-\frac{\partial}{\partial x}\otimes dx\otimes dz+\frac{\partial}{\partial y}\otimes dz\otimes dy-\frac{\partial}{\partial y}\otimes dy\otimes dz
\]
satisfies $\left<dz,T_N(X,Y)\right>=0$ for all $X,Y\in \mathfrak{X}(\mathbb R^3)$. Moreover,  $d_N(dz)=-d(i_Ndz)=0$ since $i_Ndz=2dy$. 
Therefore the pair $(N,dz)$ is a Haantjes structure on $\mathbb R^3$. 
However, a simple computation shows that a non-trivial Poisson bivector $\pi$ such that 
$\pi^{\sharp}\circ N^*=N\circ \pi^{\sharp}$ does not exist.
%\begin{equation}
%\begin{aligned}
%&\begin{pmatrix}
%0 & a & b \\ -a & 0 & c \\ -b & -c &0
%\end{pmatrix}
%\begin{pmatrix}
%0 & 2 & 0 \\ 0 & 0 & 2 \\ -\frac{y}{2} & -z & 0
%\end{pmatrix}
%=
%\begin{pmatrix}
%0 & 0 & -\frac{y}{2} \\ 2 & 0 & -z \\ 0 & 2 & 0
%\end{pmatrix}
%\begin{pmatrix}
%0 & a & b \\ -a & 0 & c \\ -b & -c & 0
%\end{pmatrix}\\
%&\Rightarrow
%\begin{pmatrix}
%-\frac{by}{2} & -bz & 2a \\ -\frac{cy}{2} & -2a-cz & 0 \\ 0 & -2b & -2c 
%\end{pmatrix}
%=
%\begin{pmatrix}
%\frac{by}{2} & \frac{cy}{2} & 0 \\ bz & 2a+cz & 2b \\ -2a & 0 & 2c
%\end{pmatrix}
%\Rightarrow a=b=c=0,
%\end{aligned}
%\end{equation}
This entails that the Haantjes structure $(N,dz)$ cannot be derived from a PqN structure.
\end{remark}

\subsection{Generalized Lenard-Magri chains}

Haantjes manifolds were introduced by Magri in \cite{MagriHaa} with the aim of generalizing the notion of Lenard-Magri chain, see Definition \ref{def:gen-LM} below. In this more general geometrical setting, the sequence of Nijenhuis torsionless tensors 
$\{N^i\}%_{i\geq 0}
$, appearing in the PN case, is traded for a family of $(1,1)$ tensors with vanishing Haantjes torsion. More precisely,


\begin{defi}[Generalized Lenard-Magri chains]\label{def:genchain} 
\label{def:gen-LM}
A generalized Lenard-Magri chain of length $p$ on a Haantjes manifold $M$ with Haantjes structure $(N,\theta)$ is a family  
$\mathcal N=\{N_0,N_1,\dots,N_{p-1}\}%_{i\geq 0}^{p-1}
$ of $(1,1)$ tensor fields satisfying, for all $i,j= 0,\dots,p-1$:
\begin{itemize}
\item[(C0)] $N_0=I$ and $N_1=N$;
\item[(C1)] $H_{N_i}=0$; 
\item[(C2)] $N_i\circ N_j=N_j\circ N_i$; 
\item[(C3)] $d_{N_i}\theta=0$; 
\item[(C4)] $\langle\theta,T_{N_i}(X,Y)\rangle=0$ 
for all $X,Y\in\mathfrak X(M)$. 
\end{itemize}
\end{defi}

The relevance of these conditions relies on the remarkable observation that, on a Haantjes manifold endowed with a generalized Lenard-Magri chain, the 1-forms $\theta_{ij}=N_iN_j\theta$ are closed --- see Proposition 1 in \cite{MagriHaa}.

\begin{remark}
The name generalized Lenard-Magri chain is not universally adopted. In particular, in \cite{MagriHaa} the structure introduced in Definition \ref{def:genchain} is simply called a Lenard chain. 
\end{remark} 

\begin{example}
Considering the same (1,1) tensor field $N$ as in Remark \ref{rem:Eber}, we notice that
\[
N^2=4\frac{\partial}{\partial z}\otimes dx-\left(y\frac{\partial}{\partial x}+2z\frac{\partial}{\partial y}\right)\otimes dy-\left(y\frac{\partial}{\partial y}+2z\frac{\partial}{\partial z}\right)\otimes dz
\]
has Nijenhuis torsion
\[
T_{N^2}=-8\frac{\partial}{\partial y}\otimes dx\otimes dy+8\frac{\partial}{\partial y}\otimes dy\otimes dx-8\frac{\partial}{\partial z}\otimes dx\otimes dz+8\frac{\partial}{\partial z}\otimes dz\otimes dx.
\]
Thus $\mathcal N=\{I,N,N^2\}$ is not a generalized Lenard-Magri chain since $\left<dz,T_{N^2}(X,Y)\right>$ does not vanish for all $X,Y\in\mathfrak{X}(\mathbb R^3)$.
\end{example}

Let $(N,\theta)$ be the Haantjes structure, with $\theta$ given by \eqref{eq: 1-form theta}, 
derived from a PqN-structure $(\pi,N,\phi)$ on an oriented three-dimensional manifold $M$ 
with volume form $\mathcal{V}$, see Theorem \ref{theo:haan}. 
\begin{theorem}
\label{theo:Len-Ma-Ge} 
The family $\mathcal N=\{N^0,N^1,\dots,N^{p-1}\}$
is a generalized Lenard-Magri chain associated to the Haantjes structure $(N,\theta)$.
\end{theorem} 
\begin{proof}
To prove the statement, one needs to check that conditions (C1), (C3) and (C4) are satisfied. To this end, first observe that Lemma \ref{lem:Bogo}-1 entails that  $H_{N^k}=0$ for every $k\geq 2$, since $H_N=0$.  
As far as (C3) is concerned, 
we notice that, for every $X\in\mathfrak X(M)$,
\[
\langle i_{N^k}\theta,X\rangle=\langle\theta,N^kX\rangle=\langle i_{N}\theta,N^{k-1}X\rangle\stackrel{\eqref{eq:iNxi}}{=}(\lambda+\langle\xi,Z\rangle)\langle i_{N^{k-1}}\theta,X\rangle,
\]
which entails
\begin{equation}
i_{N^k}\theta=(\lambda+\langle\xi,Z\rangle)^k\theta.\label{eq:iNxi1}
\end{equation}
Since $\theta$ is closed 
\[
d_{N^k}\theta=-di_{N^k}\theta=-d(\lambda+\langle\xi,Z\rangle)^k\theta=-k(\lambda+\langle\xi,Z\rangle)^{k-1}d(\lambda+\langle\xi,Z\rangle)\wedge\theta\stackrel{\eqref{eq: 3D PqN conditions}}{=}0.
\]
To prove that (C4) holds, we simply observe that $\left<\theta,T_{N^k}(X,Y)\right>$ vanishes because of (\ref{eq: T_N 3D PqN-k}) and the definition of $\theta$.
\end{proof}


\subsection{Comparison with the work of Tempesta and Tondo}\label{ss:tempton}
We will now make a few comments to link the theory presented in the previous sections with the recent works of Tempesta and Tondo, see \cite{tempestatondoben, tempestatondo, tempestatondohigher,tempestatondoclass} and \cite{TondoLag}. We start recalling %that  
the notion of Haantjes algebra, %was 
introduced in \cite{tempestatondo}, %see 
Definition 36. %in loc. cit.. 
This consists of a set of $(1,1)$ tensor fields $\{N_i\}_{i\in I}$ on a manifold $M$ 
%, having zero Haantjes torsion and 
satisfying the following two properties:
\begin{itemize}
\item $H_{fN_i+gN_j}=0$, for all $i,j\in I$ and $f,g\in C^\infty(M)$;
\item $H_{N_iN_j}=0=H_{N_jN_i}$, for all $i,j\in I$.
\end{itemize}
%In particular, $H_{N_i}=0$ for all $i\in I$.  \blu{[We should change the following sentence.]}
%Note that the first property says that the $N_i$ %s 
%generate a $C^\infty(M)$-module and that the two properties, if taken together, say that the $N_i$ %s 
%generate a ring.

A Haantjes algebra will be called of rank $m$ if the corresponding $C^\infty(M)$-module has rank equal to $m$. The Haantjes algebra will be called abelian if the operators $N_i$ %s 
pairwise commute. An important class of abelian Haantjes algebras are the cyclic ones, i.e., the Haantjes algebras generated by a given Haantjes operator. A cyclic Haantjes algebra has rank $m$ if the minimal polynomial of its generator has degree $m$. 

After these preliminary comments, we easily obtain %\blu{[using Theorem \ref{theo:Len-Ma-Ge} ? Why rank-two?] }
\begin{prop}\label{pro:TTpqn}
If $(\pi,N,\phi)$ is a PqN structure on an oriented three-dimensional manifold and $\pi$ never vanishes, 
then $N$ generates a rank-two %, 
cyclic Haantjes algebra.
\end{prop}

\begin{proof}
As we already noticed in Theorem \ref{theo:haan}, $N$ is a Haantjes operator.  Then $\lbrace N^k\rbrace_{k\geq 0}$ is a cyclic Haantjes algebra as a consequence of Lemma \ref{lem:Bogo}-1. Moreover, the rank of such algebra is two since the minimal polynomial of $N=\lambda I+Z\otimes\xi$ is  $p(t)=t^2-(\left<\xi,Z\right>+2\lambda)t+\lambda^2+\left<\xi,Z\right>\lambda$.
\end{proof}

Given a Haantjes algebra of rank $m$ and a set $N_1,\dots,N_m$ of its generators (as a $C^\infty(M)$-module), a closed 1-form $\alpha$ generates a so called \emph{Haantjes chain} $(N_1^\ast\alpha,\dots,N_m^\ast\alpha)$ of length $m$ if
%, for all $k=1,\dots,m$,
$
d_{N_k}\alpha=0,
$
i.e., if $N^\ast_k\alpha=i_{N_k}\alpha$ is closed for all $k=1,\dots,m$, and the forms $N_k^\ast\alpha$, with $k=1,\dots,m$, are linearly independent, see Definition 15 in \cite{tempestatondoclass}.

If $(\pi,N,\phi)$ is as in Proposition \ref{pro:TTpqn}, where $N=\lambda I+Z\otimes\xi$ 
and $\theta$ given as in \eqref{eq: 1-form theta},
%$d\xi=0$
then the 1-forms 
%$\xi_k={N^k}^\ast\xi=i_{N^k}\xi$, 
$\theta_k={N^k}^{\ast}\theta=i_{N^k}\theta$, 
see \eqref{eq:iNxi1}, are closed but, clearly, they are not independent (since %$\xi$
$\theta$ is a common eigenvector of the ${N^k}^\ast$) and for this reason they fail to form a Haantjes chain.

%On the other hand, under the same assumptions on the PqN structure made in Proposition \ref{pro:TTpqn}, every closed 1-form $\alpha$ such that 
%\begin{enumerate}
%\item $\xi$ and $\alpha$ are linearly independent,
%\item $\langle\alpha,Z\rangle$ is not identically zero,
%\end{enumerate}
%generates a Haantjes chain on the dense open set where $\langle\alpha,Z\rangle\neq 0$.
Finally, it is worth mentioning that every PqN structure satisfying the hypothesis of Proposition \ref{pro:TTpqn} is a \emph{Poisson-Haantjes structure} in the sense of Definition 6 in \cite{TondoLag}.



\section{Loose ends}\label{sec:Appendix}

We divide this section in three parts. The first one is devoted to the proof of \eqref{eq:dNphi},
%\textcolor{blue}{\eqref{eq:phi}},
the second one to an alternative proof of 
Proposition \ref{thm:inv3}  and the last one is aimed to collect some more information about the Haantjes torsion.

\subsection{Proof of the identity \eqref{eq:dNphi}}\label{subsec:app1}
First notice that, given $\pi=\frac{\partial}{\partial x}\wedge\frac{\partial}{\partial y}$  and any closed 2-form 
$\Omega=\Omega_{12}dx\wedge dy+\Omega_{13}dx\wedge dz+\Omega_{23}dy\wedge dz$, we have
\[
\frac{1}{2}\left[\Omega,\Omega\right]_{\pi}=\left(\Omega_{23}\frac{\partial\Omega_{12}}{\partial x}-\Omega_{13}\frac{\partial\Omega_{12}}{\partial y}+\Omega_{12}\frac{\partial\Omega_{12}}{\partial z}\right)\mathcal V.
\]
In addition, for $N_1=\widetilde{Z}\otimes dz$ as considered in %Example \ref{exa: PqN-R3 deformed}
Proposition \ref{prop:local}, one has that
\begin{equation}\label{eq: d_N1Omega} %%%%%%%
d_{N_1}\Omega=-di_{N_1}\Omega=-d\left(i_{N_1}(\Omega_{12}dx\wedge dy)+i_{N_1}(\Omega_{13}dx\wedge dz)+i_{N_1}(\Omega_{23}dy\wedge dz)\right).
\end{equation}
%Since $\Omega=\Omega_{ij}dx_i\wedge dx_j$ is closed, then $d_{N_1}\Omega=-di_{N_1}\Omega$.
%\[
%d_{N_1}\Omega=(i_{N_1}\circ d-d\circ i_{N_1})\Omega=-di_{N_1}\Omega.
%\]
%Notice that
%\[
%i_{N_1}\Omega=i_{N_1}(\Omega_{12}dx\wedge dy)+i_{N_1}(\Omega_{13}dx\wedge dz)+i_{N_1}(\Omega_{23}dy\wedge dz).
%\]
%Computing each of the summands of the previous sum we have
We observe that
\begin{align*}
i_{N_1}(\Omega_{12}dx\wedge dy)(X,Y)%&=(\Omega_{12}dx\wedge dy)(N_1X,Y)+(\Omega_{12}dx\wedge dy)(X,N_1Y)\\
 &=(\Omega_{12}dx\wedge dy)\left(g(z)dz(X)\frac{\partial}{\partial z},Y\right)+(\Omega_{12}dx\wedge dy)\left(X,g(z)dz(Y)\frac{\partial}{\partial z}\right)=0,
\end{align*}
\begin{align*}
i_{N_1}(\Omega_{13}dx\wedge dz)(X,Y)%&=(\Omega_{13}dx\wedge dz)(N_1X,Y)+(\Omega_{13}dx\wedge dz)(X,N_1Y)\\
&=(\Omega_{13}dx\wedge dz)\left(g(z)dz(X)\frac{\partial}{\partial z},Y\right)+(\Omega_{13}dx\wedge dz)\left(X,g(z)dz(Y)\frac{\partial}{\partial z}\right)\\
%&=\cancel{(\Omega_{13}dx)(g(z)dz(X)\partial_z)}dz(Y)-(\Omega_{13}dx)(Y)dz(g(z)dz(X)\partial_z)+\\
%\;\;\;\;&+(\Omega_{13}dx)(X)dz(g(z)dz(Y)\partial_z)-\cancel{(\Omega_{13}dx)(g(z)dz(Y)\partial_z)}dz(X)\\
&=g(z)\Omega_{13}dx(X)dz(Y)-g(z)\Omega_{13}dx(Y)dz(X)\\
&=g(z)(\Omega_{13}dx\wedge dz)(X,Y),\\
%\end{align*}
%and, finally,
%\begin{align*}
i_{N_1}(\Omega_{23}dy\wedge dz)(X,Y)%&=(\Omega_{23}dy\wedge dz)(N_1X,Y)+(\Omega_{23}dy\wedge dz)(X,N_1Y)\\
&=(\Omega_{23}dy\wedge dz)\left(g(z)dz(X)\frac{\partial}{\partial z},Y\right)+(\Omega_{23}dy\wedge dz)\left(X,g(z)dz(Y)\frac{\partial}{\partial z}\right)\\
%&=\cancel{(\Omega_{23}dy)(g(z)dz(X)\partial_z)}dz(Y)-(\Omega_{23}dy)(Y)dz(g(z)dz(X)\partial_z)+\\
%\;\;\;\;&+(\Omega_{23}dy)(X)dz(g(z)dz(Y)\partial_z)-\cancel{(\Omega_{23}dy)(g(z)dz(Y)\partial_z)}dz(X)\\
&=g(z)\Omega_{23}dy(X)dz(Y)-g(z)\Omega_{23}dy(Y)dz(X)\\
&=g(z)(\Omega_{23}dy\wedge dz)(X,Y).
\end{align*}
%Therefore
%Using the three previous equations and put them in \eqref{eq: d_N1Omega} gives us
By substituting the three preceding equations into \eqref{eq: d_N1Omega}, we obtain
%\begin{align*}
%i_{N_1}\Omega &=i_{N_1}(\Omega_{12}dx\wedge dy)+i_{N_1}(\Omega_{13}dx\wedge dz)+i_{N_1}(\Omega_{23}dy\wedge dz)\\
%&=g(z)(\Omega_{13}dx\wedge dz+\Omega_{23}dy\wedge dz)\\
%&=g(z)\Omega-g(z)\Omega_{12}dx\wedge dy.
%\end{align*}
%and
%Then we have
\begin{align*}
d(i_{N_1}\Omega)&=d(g(z)\Omega-g(z)\Omega_{12}dx\wedge dy)=dg(z)\wedge \Omega-\frac{\partial \left(g(z)\Omega_{12}\right)}{\partial z}\mathcal V\\
&=g(z)'dz\wedge (\Omega_{12}dx\wedge dy+\Omega_{13}dx\wedge dz+\Omega_{23}dy\wedge dz)-\left(g(z)'\Omega_{12}+g(z)\frac{\partial \Omega_{12}}{\partial z}\right)\mathcal V\\
%&=g(z)'\Omega_{12}\mathcal V-g(z)'\Omega_{12}\mathcal V-g(z)(\Omega_{12})_z\mathcal V\\
&=-g(z)\frac{\partial \Omega_{12}}{\partial z}\mathcal V.
\end{align*}
Thus if $\Omega$ is a closed 2-form on $\mathbb{R}^3$, then 
%\[
%d_{N_1}\Omega=-d(i_{N_1}\Omega)=g(z)%\cdot
%(\Omega_{12})_z\mathcal V.
%\]
%Since
%Finally we have
\begin{equation}\label{eq: deformed_part}
d_{N_1}\Omega+\frac{1}{2}\left[\Omega,\Omega\right]_{\pi}=\left(-g(z)\frac{\partial \Omega_{12}}{\partial z}+\Omega_{23}\frac{\partial\Omega_{12}}{\partial x}-\Omega_{13}\frac{\partial\Omega_{12}}{\partial y}+\Omega_{12}\frac{\partial\Omega_{12}}{\partial z}\right)\mathcal V.
\end{equation}
In particular, the 2-form $\Omega$ given in %Example \ref{exa: PqN-R3 deformed} 
Proposition \ref{prop:local} is closed and from \eqref{eq: deformed_part} we derive 
\[
d_{N_1}\Omega+\frac{1}{2}\left[\Omega,\Omega\right]_{\pi}=(g(z)\lambda_z-a\lambda_x-b\lambda_y+\lambda\lambda_z)\mathcal V=-Z(\lambda)\mathcal V=\phi,
\]
since $\Omega_{12}=-\lambda,\Omega_{13}=-b,\Omega_{23}=a$ and $\lambda+c=g(z)$.

%\begin{eqnarray}%\label{eq: particular [Omega,Omega]_{pi}}
%\frac{1}{2}\left[\Omega,\Omega\right]_{\pi}%&=& \left(a(-\lambda)_x+b(-\lambda)_y-\lambda(-\lambda)_z\right)\mathcal V
%\nonumber\\
%&=&-(a\lambda_x+b\lambda_y-(g(z)-c)\lambda_z)\mathcal V=-Z(\lambda)\mathcal V+g(z)\lambda_z\mathcal V.\label{eq: particular [Omega,Omega]_{pi}}  
%\end{eqnarray}


\subsection{An alternative proof of Proposition \ref{thm:inv3}}\label{ss:proofth} In this subsection we present an alternative proof of Proposition \ref{thm:inv3} which does not use local coordinates. In spite of being more involved, it is interesting per-se and we think it could play a role in a possible extension of the results of this paper to higher (greater than 3) dimension.

The proof is based on a closed formula for the 1-forms $\phi_k$, see \eqref{eq:genrec}, which is contained in the following
\begin{lem}\label{lem:phik}
	On an oriented three-dimensional PqN manifold, where $N=\lambda I+Z\otimes\xi$, we have that 
	\begin{equation}
		\phi_k=Z(\lambda)\lambda^k\xi\qquad\mbox{for all $k\geq 0$.}
		\label{eq:phik}
	\end{equation}
\end{lem}
\begin{proof}
	The first step is to prove %(see Section \ref{subsec:app2}) 
	that 
	\begin{equation}
		N^k=\lambda^kI+f_k \,Z\otimes\xi,\label{eq:Nalk}
	\end{equation}
	where $f_k=\sum_{l=0}^{k-1}{k\choose l}\lambda^l\langle\xi,Z\rangle^{k-l-1}$. 
The proof is done by induction, noticing that for $k=1$ one has $N=\lambda I+Z\otimes \xi$, which is the content of Theorem \ref{theo:charN}.
Suppose now that \eqref{eq:Nalk} holds for $k$ and let us prove it for $k+1$. We compute
\begin{equation}
\label{eq:s1}
\begin{aligned}
N^{k+1}(X)&=(\lambda I+Z\otimes\xi)\Big(\lambda^kX+\sum_{l=0}^{k-1}{k\choose l}\lambda^l\langle\xi,Z\rangle^{k-l-1}\langle\xi,X\rangle Z\Big)
%\nonumber
\\
&=\lambda^{k+1}X+\underbrace{\sum_{l=0}^{k-1}{k\choose l}\lambda^{l+1}\langle\xi,Z\rangle^{k-l-1}\langle\xi,X\rangle Z}_{(1)}+\underbrace{\sum_{l=0}^{k-1}{k\choose l}\lambda^{l}\langle\xi,Z\rangle^{k-l}\langle\xi,X\rangle Z+\lambda^k\langle\xi,X\rangle Z}_{(2)}.
\end{aligned}
\end{equation}
Now we note that in \eqref{eq:s1} one has
\begin{eqnarray}
&(1)&\stackrel{s=l+1}{=}\sum_{s=1}^{k}{k\choose s-1}\lambda^{s}\langle\xi,Z\rangle^{k-s}\langle\xi,X\rangle Z,\label{eq:s2}\\
&(2)&=\sum_{l=0}^k{k\choose l}\lambda^l\langle\xi,Z\rangle^{k-l}\langle\xi,X\rangle Z=\sum_{l=1}^k{k\choose l}\lambda^l\langle\xi,Z\rangle^{k-l}\langle\xi,X\rangle Z+\langle\xi,Z\rangle^k\langle\xi,Z\rangle Z.\label{eq:s3}
\end{eqnarray}
By renaming $l$ the running variable $s$ in \eqref{eq:s2}, plugging \eqref{eq:s2} and \eqref{eq:s3} in \eqref{eq:s1}, 
and using the identity ${k\choose l-1}+{k\choose l}={k+1\choose l}$, one obtains 
\begin{eqnarray*}
%&&
N^{k+1}(X)%\\
&=&\lambda^{k+1}X+\sum_{l=1}^{k}\Big[{k\choose l-1}+{k\choose l}\Big]\lambda^{l}\langle\xi,Z\rangle^{k-l}\langle\xi,X\rangle Z+\langle\xi,Z\rangle^k\langle\xi,Z\rangle Z\\
&%\stackrel{{k\choose l-1}+{k\choose l}={k+1\choose l}}{=}
=&\lambda^{k+1}X+\sum_{l=1}^{k}{k+1\choose l}\lambda^{l}\langle\xi,Z\rangle^{k-l}\langle\xi,X\rangle Z+\langle\xi,Z\rangle^k\langle\xi,Z\rangle Z\\
&=&\lambda^{k+1}X+\sum_{l=0}^{k}{k+1\choose l}\lambda^{l}\langle\xi,Z\rangle^{k-l}\langle\xi,X\rangle Z,
\end{eqnarray*}
which, being $X$ arbitrary, entails the identity \eqref{eq:Nalk}.  

The second step starts with the remark that, since
	\[
	T_N(X,Y)\stackrel{\eqref{eq: T_N 3D PqN}}{=}Z(\lambda)i_\xi(X\wedge Y)=Z(\lambda)\langle\xi,X\rangle Y-Z(\lambda)\langle\xi,Y\rangle X,
	\]
	one has that
	\begin{equation}
		i_XT_N=Z(\lambda)(\langle\xi,X\rangle I-X\otimes\xi).\label{eq:iXTN}
	\end{equation}
	For every $k\geq 0$ and $X\in\mathfrak X(M)$, we can now compute
	\[
	\begin{split}
		\langle\phi_k,X\rangle&\stackrel{%\eqref{eq:phi}
			\eqref{eq:phis}}{=}\frac12\text{Tr}(N^ki_XT_N)\stackrel{\eqref{eq:iXTN}}
		{=}\frac12\text{Tr}\big(N^k(Z(\lambda)\langle\xi,X\rangle I-Z(\lambda)X\otimes\xi)\big)\\
		&\stackrel{\eqref{eq:Nalk}}{=}\frac12\text{Tr}\Big[\Big(\lambda^kI+\sum_{l=0}^{k-1}{k\choose l}\lambda^l\langle\xi,Z\rangle^{k-l-1}Z\otimes\xi \Big)\Big(Z(\lambda)\langle\xi,X\rangle I-Z(\lambda)X\otimes\xi)\Big)\Big]\\
		&=\frac12\text{Tr}\Big[Z(\lambda)\lambda^k\langle\xi,X\rangle I-Z(\lambda)\lambda^kX\otimes\xi+Z(\lambda)\langle\xi,X\rangle\sum_{l=0}^{k-1}{k\choose l}\lambda^l\langle\xi,Z\rangle^{k-l-1}Z\otimes\xi\\
		&-Z(\lambda)\sum_{l=1}^{k-1}{k\choose l}\lambda^l\langle\xi,Z\rangle^{k-l-1}\langle\xi,X\rangle Z\otimes\xi\Big],
	\end{split}
	\]
	where, in the last sum, $\langle\xi,X\rangle Z\otimes\xi$ comes from the  computation
	\[
	\left((Z\otimes\xi)\circ (X\otimes\xi)\right)(Y)=\langle\xi,Y\rangle(Z\otimes\xi)(X)=\langle\xi,Y\rangle\langle\xi,X\rangle Z%,
	\qquad\forall Y\in\mathfrak X(M),
	\]
	which yields the identity $(Z\otimes\xi)\circ(X\otimes\xi)=\langle\xi,X\rangle Z\otimes\xi$. Now, recalling that $\text{Tr}(Z\otimes\xi)=\langle\xi,Z\rangle$ and $\text{Tr}(I)=3$, the last term of the previous chain of identities becomes
	\[
	\begin{split}
		&\frac12\left(3Z(\lambda)\lambda^k\langle\xi,X\rangle-\lambda^kZ(\lambda)\langle\xi,X\rangle+\cancel{Z(\lambda)\langle\xi,X\rangle\sum_{l=0}^{k-1}{k\choose l}\lambda^l\langle\xi,Z\rangle^{k-l}}\cancel{-Z(\lambda)\langle\xi,X\rangle\sum_{l=0}^{k-1}{k\choose l}\lambda^l\langle\xi,Z\rangle^{k-l}}
		\right)\\
		&=Z(\lambda)\lambda^k\langle\xi,X\rangle,
	\end{split}
	\]
	which entails $\phi_k=Z(\lambda)\lambda^k\xi$.
\end{proof}
With this result at hand, we can go back to the
\begin{proof}(of %Theorem
	Proposition \ref{thm:inv3}) We want to prove that $\{I_k,I_j\}=0$ for all $k,j$ with $k>j$. 
	Since the 1-forms $\phi_k$ are multiple of $\xi$, see \eqref{eq:phik}, and $\pi^\sharp\xi=0$, see (\ref{prop:xi2}),
	the right-hand side of \eqref{eq:recursionphi} is %(identically) 
	zero. In other words, that identity collapses to
	\[
	\{I_k,I_j\}=\{I_{k-1},I_{j+1}\}.
	\] 
	If $k-1=j+1$ we are done, otherwise one can apply again this argument, which, after a finite number of steps, will entail the thesis.
\end{proof}

\begin{remark} 
	\label{rem:torsion-of-k-power}
	Using (\ref{eq: T_N computation}) and (\ref{eq:Nalk}), one obtains the following 
	generalization of (\ref{eq: T_N 3D PqN}), giving the Nijenhuis torsion of the $k$-th power of the tensor 
	$N=\lambda I+Z\otimes\xi$:
	\begin{equation}
		T_{N^k}(X,Y)=f_k Z(\lambda^k)i_\xi(X\wedge Y).\label{eq: T_N 3D PqN-k}
	\end{equation}
	This could also be deduced from the general expression in Remark 3 of \cite{BogoTens}.
\end{remark}

We close this subsection with the following remark, where we compare the result just obtained with Theorem \ref{theo:inv}, giving sufficient conditions for the involutivity of a (general) PqN manifold.
\begin{remark} \label{rem:FMP2024}
	Applying $\pi^\sharp$ to both sides of \eqref{eq:genrec} and using again (\ref{prop:xi2}) and Lemma \ref{lem:phik}, one arrives to the identity $X_{k+1}=N X_k$, which holds for all $k\geq 1$ and which, if iterated, gives $X_{k+1}=N^{k}X_1$ for all $k\geq 1$.
	This observation implies that the vector fields $Y_k$ appearing in item (b) of Theorem \ref{theo:inv} are identically zero. 
	On the other hand, one can find a 2-form $\Omega$ such that $\phi=-Z(\lambda)\mathcal V=-2dI_1\wedge\Omega$, see for example \cite{Pham}, Chapter 3, Section 1.1. From these observations we can deduce that both conditions of Theorem \ref{theo:inv} are satisfied, and we obtain another proof of %Theorem
	Proposition \ref{thm:inv3}.  
\end{remark}

\subsection{Recollection on the Haantjes torsion}\label{ss:haantjes}

To put the definition \eqref{eq:haant} of the Haantjes torsion in context, we need to introduce a few notions, which we enclose in the following 

\begin{defi}[See \cite{tempestatondo} and \cite{BenentiChanuRastelli}]
\begin{itemize}
\item[]
\item A frame over the open set $U\subset M$ is a collection ${\mathcal X}:=\{X_1,\dots,X_n\}$ of $n=\dim M$ vector fields which generate the tangent space at each point of $U$. 
\item Two frames $\mathcal X:=\{X_1,\dots,X_n\}$ and $\mathcal Y:=\{Y_1,\dots,Y_n\}$ over $U$ are called equivalent if there exists a collection of never vanishing smooth functions $f_1,\dots,f_n$ over $U$ such that $X_i=f_i Y_i$, for all $i=1,\dots,n$.
\item A %reference 
frame $\mathcal X=\{X_1,\dots,X_n\}$ over $U$ is called integrable if for all $p\in U$ there exists a chart around $p$, 
say $(V,x_1,\dots,x_n)$, such that the restriction of $\mathcal X$ to $V$ is equivalent to the frame $\{\frac{\partial}{\partial x_1},\dots,\frac{\partial}{\partial x_n}\}$.
\item A tensor field $N$ of type $(1,1)$ is called semi-simple if each point $p\in M$ has a neighborhood $U$ where one can find a %\emph{reference 
frame %} 
formed by eigenvectors of $N$. Such a frame is said to be an eigenframe for $N$.
\end{itemize}
\end{defi}

It is also worth recalling %the following result
\begin{prop}%[See Proposition 19 in \cite{tempestatondo}] 
A reference frame $\mathcal X=\{X_1,\dots,X_n\}$ on $U$ is integrable if and only if one of the two following equivalent conditions is verified:
\begin{enumerate}
\item For every $i,j\in\{1,\dots,n\}$, the distribution generated by $\{X_i,X_j\}$ is Frobenius integrable.
\item The distribution generated by $\mathcal X\backslash\{X_i\}$ is Frobenius integrable.
\end{enumerate}
\end{prop}

The integrability of eigenframes for a given $(1,1)$ tensor field is controlled by the Haantjes torsion of $N$. 
More precisely, the vanishing of the Haantjes torsion of $N$ is a necessary and sufficient condition for the integrability of its eigenframes, see Theorem 25 in \cite{tempestatondo} and references therein. 

\section{Final comments and outlook}\label{sec:comfin} PN and bi-Hamiltonian structures has been proven crucial in the analysis of the complete integrability of classical Hamiltonian systems. In the series of papers \cite{FMOP2020,FMP2023,FMP2024}, it has been observed that under some stringent assumptions, one could use the notion of a PqN structure to describe the complete integrability of some classical and important integrable system from a geometrical point of view. To this end, the notion of involutive PqN structure was introduced and the one of a deformation of a PqN structure was analyzed.

In this paper first we characterize the PqN structures defined on oriented three dimensional manifolds whose underlying Poisson tensor never vanishes, see Theorem \ref{theo:charN}. Then we show that every such PqN structure is involutive, see Proposition \ref{thm:inv3} and that, \emph{locally}, it comes from a deformation of a PN one, see Proposition \ref{prop:local}. Theorem \ref{theo:haan} links this class of PqN structures with the Haantjes ones, while Theorem \ref{theo:Len-Ma-Ge} is a first attempt to build a bridge between the PqN structures and the so called generalized Lenard-Magri chains.

We believe that the results obtained in this work, although interesting in themselves, are preliminary to future research.
A few directions we would like to explore are the following ones.

First, we would like to investigate the possibility to extend to higher dimension the results enclosed in Section \ref{sec:3}. More precisely, as a preliminary checks, one could try to see if $\mathbb R^{2n+1}$, $n>1$, carries a PqN structure whose (1,1) tensor has the analogue form of \eqref{eq:thanksRef} and whose Poisson tensor, in Darboux coordinates $(x^1,\dots,x^n,y_1,\dots,y_n,z)$, is of the type $\pi=\sum_{i=1}^n\frac{\partial}{\partial x^i}\wedge\frac{\partial}{\partial y_i}$. If such structures exist, one could ask if for them the results in Propositions \ref{prop:local} and \ref{thm:inv3} hold true. 

Another direction we would like to pursue is to investigate under what assumptions a PqN structure in dimension greater than 3 is \emph{Haantjes torsion free}, see Theorem \ref{theo:haan}. We think that anything in this direction could be relevant to connect the theory of PqN structures to the one of dynamical systems of hydrodynamical type, see for example \cite{DN1,DN2}. It is already an interesting problem to look for systems of hydrodynamical type which can be defined from the PqN structures described in this paper, see \cite{FMS}. 

Finally, it could be interesting to see if the result  in Subsection \ref{ss:tempton}
%\ref{ss:tempton} 
could be extended to higher dimension.

%\blu{[Is semisimplicity needed for this result? If not, why do we recall the definition of semisimplicity?]  Possible asnwer: Such definition is related with ``eigenframe'' which is mentioned in the last paragraph.}

%\bigskip\bigskip

%\blu{\todo{DA FARE SOLO PRIMA DI SPEDIRE IL LAVORO: controllare se le references qui sotto sono tutte citate.}}

\thebibliography{99}

\bibitem{Antunes2008}
Antunes, P., {\it Poisson quasi-Nijenhuis structures with background}, Lett. Math. Phys. {\bf 86} (2008), 33--45.

\bibitem{BenentiChanuRastelli} Benenti, S., Chanu, C., Rastelli, G., {\it Remarks on the connection between the additive separation of the Hamilton-Jacobi equation and the multiplicative separation of the Schr\"odinger equation. I. The completeness and Robertson conditions}, J. Math. Phys. {\bf 43} (2002), 5183--5222.
%\bibitem{Bogo96-180} Bogoyavlenskij, O.I., {\it Theory of Tensor Invariants of Integrable Hamiltonian Systems. I. Incompatible Poisson Structures}, 
%Commun. Math. Phys. {\bf 180} (1996), 529--586.

\bibitem{Bogo96-182}  Bogoyavlenskij, O.I., {\it Necessary Conditions for Existence of Non-Degenerate Hamiltonian Structures}, Commun. Math. Phys. {\bf 182} (1996), 253--290.

\bibitem{BogoTens} Bogoyavlenskij, O.I., {\it Algebraic identities for the Nijenhuis tensor}, Diff. Geom. Appl. {\bf 24} (2006) 447--457.

%\bibitem{BKM2022}
%{Bolsinov, A.V., Konyaev, A.Yu., Matveev, V.S.}, {\it Nijenhuis geometry}, {Adv. Math.} {\bf 394} (2022), 52 pages.

\bibitem{Bonechi} Bonechi, F., {\it  Multiplicative integrable models from Poisson-Nijenhuis structures}, in: {\rm From Poisson brackets to universal quantum symmetries}, Banach Center Publ. {\bf 106}, Warsaw, 2015, pp.\ 19--33.

%\bibitem{BursztynDrummond2019}  Bursztyn, H., Drummond, T.,
%{\it Lie theory of multiplicative tensors}, Math.\ Ann.\ {\bf 375} (2019), 1489--1554.

\bibitem{BursztynDrummondNetto2021} Bursztyn, H., Drummond, T., Netto, C.,
{\it Dirac structures and Nijenhuis operators}, Math. Z. {\bf 302} (2022), 875--915.

%https://doi.org/10.48550/arXiv.2109.06330

\bibitem{Cheng-Sheng} Chen, B., Sheng, Y., {\it Poisson-Nijenhuis structures on oriented 3D-manifolds}, 
Rep. Math. Phys. {\bf 61} (2008), 361--380. 

\bibitem{FMP2024}  Chu\~no Vizarreta, E., Falqui, G., Mencattini, I., Pedroni, M.,  {\it Poisson quasi-Nijenhuis manifolds, closed Toda lattices, and generalized recursion relations}, Lett. Math. Phys. {\bf 115}, no. 84 (2025), 22 pages. 

\bibitem{FMP2026}  Chu\~no Vizarreta, E., Falqui, G.; Mencattini, I., Pedroni,  M., {\it An involutivity theorem for a class of Poisson quasi-Nijenhuis manifolds}, arXiv:2603.06532. %, 22 pages. 

\bibitem{C-NdC-2010} Cordeiro, F., Nunes da Costa, J.M.,
{\it Reduction and construction of Poisson quasi-Nijenhuis manifolds with background}, 
Int. J. Geom. Methods Mod. Phys. {\bf 7} (2010), 539--564.

%\bibitem{Damianou99} Damianou, P.A., {\it Multiple Hamiltonian structures for Toda systems of type A-B-C}, Regul. Chaotic Dyn. {\bf 5} (2000), 17--32.

%\bibitem{Damianou04rev} Damianou, P.A., {\it Multiple Hamiltonian structures of Bogoyavlensky-Toda lattices}, Rev. Math. Phys. {\bf 16} (2004), 175--241.

%\bibitem{Damianou04} Damianou, P.A., {\it On the bi-Hamiltonian structure of Bogoyavlensky-Toda lattices}, 
%Nonlinearity {\bf 17} (2004), 397--413. 

\bibitem{DO} Das, A., Okubo, S., {\it A systematic study of the Toda lattice}, Ann. Physics {\bf 190} (1989), 215--232.

\bibitem{DMP2024} do Nascimento Luiz, M., Mencattini, I., Pedroni, M.,
{\it Quasi-Lie bialgebroids, Dirac structures, and deformations of Poisson quasi-Nijenhuis manifolds\/}, 
Bull.\ Braz.\ Math.\ Soc.\ (N.S.) {\bf 55} (2024), 18 pages.

\bibitem{DN1} Dubrovin, B.A., Novikov, S.P., {\it Poisson brackets of hydrodynamic type}, Dokl. Akad. Nauk SSSR {\bf 279} (1984),  294--297.

\bibitem{DN2} Dubrovin, B.A., Novikov, S.P., {\it Hydrodynamics of weakly deformed soliton lattices. Differential geometry and Hamiltonian Theory}, Russ. Math. Surv. {\bf 44} (1989), 35--124.


%\bibitem{FMP2001} Falqui, G., Magri, F., Pedroni, M.,
%{\it Bihamiltonian geometry and separation of variables for Toda lattices\/}, J.\ Nonlinear
%Math.\ Phys.\ {\bf 8} (2001), suppl., 118--127.

%\bibitem{FMT2000} Falqui, G., Magri, F., Tondo, G.,
%{\it Reduction of bi-Hamiltonian systems and the separation of variables: an example from the Boussinesq hierarchy\/}, 
%Theoret.\ and Math.\ Phys.\ {\bf 122} (2000), 176--192.

\bibitem{FMOP2020} Falqui, G., Mencattini, I., Ortenzi, G., Pedroni, M.,
{\it Poisson quasi-Nijenhuis manifolds and the Toda system\/}, Math.\ Phys.\ Anal.\ Geom.\ {\bf 23} (2020), 17 pages.

\bibitem{FMP2023} Falqui, G., Mencattini, I., Pedroni, M.,
{\it Poisson quasi-Nijenhuis deformations of the canonical PN structure\/}, J.\ Geom.\ Phys.\  {\bf 186} (2023), 10 pages.


%\bibitem{FP03} Falqui, G., Pedroni, M., {\it Separation of variables for bi-Hamiltonian systems\/}, Math.\ Phys.\ Anal.\ Geom.\ 
%{\bf 6} (2003), 139--179.

\bibitem{FeraMar} Ferapontov, E.V., Marshall, D.G., {\it Differential-geometric approach to the integrability of hydrodynamic chains: the Haantjes tensor}, Math. Ann. {\bf 339} (2007), 61--99.

\bibitem{FMS} Ferapontov, E.V., Moro, A., Sokolov, V.V., {\it Hamiltonian Systems of Hydrodynamics Type in 2+1 Dimensions}, 
Commun. Math. Phys. {\bf 285} (2009), 31--65.  

%\bibitem{FiorenzaManetti2012} Fiorenza, D., Manetti, M.,
%{\it Formality of Koszul brackets and deformations of holomorphic Poisson manifolds},
%Homology Homotopy Appl. {\bf 14} (2012), 63--75. 

%\bibitem{ILX} Iglesias-Ponte, D., Laurent-Gengoux, C., Xu, P., {\it Universal lifting theorem and quasi-Poisson groupoids}, 
%J. Eur. Math. Soc. (JEMS) {\bf 14} (2012), 681--731. 

\bibitem{Haantjes} Haantjes, J., {\it On $X_{n-1}$-forming sets of eigenvectors}, Indag. Math. {\bf 17} (1955), 158--162.

%\bibitem{YKS96} Kosmann-Schwarzbach, Y., {\it The Lie Bialgebroid of a Poisson-Nijenhuis Manifold}, Lett. Math. Phys. {\bf 38} (1996), 421--428.

\bibitem{KM} Kosmann-Schwarzbach, Y., Magri, F., {\it Poisson-Nijenhuis structures}, Ann. Inst. Henri Poincar\'e {\bf 53} (1990), 35--81.

\bibitem{Kos-Recursion} Kosmann-Schwarzbach, Y., {\it Beyond recursion operators}, in: 
Geometric methods in physics XXXVI, Trends Math., Birkh\"auser/Springer, Cham, 2019, pp.\ 167--180.

%\bibitem{YKS-Rubtsov} Kosmann-Schwarzbach, Y., Rubtsov, V., {\it Compatible Structures on Lie Algebroids and Monge-Amp\`ere Operators}, Acta Appl. Math. {\bf 109} (2010), 101--135.

\bibitem{L-GPV} Laurent-Gengoux, C., Pichereau, A., Vanhaecke, P., {\it Poisson structures}, Grundlehren der mathematischen Wissenschaften, Volume 347, Springer-Verlag, Berlin Heidelberg, 2013.

%\bibitem{LWX} Liu, Z-J., Weinstein, A., Xu, P., {\it Manin Triples for Lie Bialgebroids}, J. Differential Geom. {\bf 45} (1997), 547--574.

%\bibitem{Magri2003} Magri, F., 
%{\it Lenard chains for classical integrable systems}, Theoret.\ and Math.\ Phys.\ {\bf 137} (2003), 1716--1722.

\bibitem{MagriHaa} Magri, F., {\it Haantjes manifolds}, Journal of Physics: Conference Series {\bf 482} (2014), 012028.

\bibitem{MagriVeselov} Magri, F., {\it Haantjes manifolds and Veselov systems}, Theoret.\ and Math.\ Phys.\ {\bf 189} (2016), 1486--1499.

\bibitem{MagriSymm} Magri, F., {\it  Haantjes manifolds with symmetries}, Theoret.\ and Math.\ Phys.\ {\bf 196} (2018), 1217--1229.

\bibitem{Magri-Morosi} Magri, F., Morosi, C., {\it A geometrical characterization of integrable Hamiltonian systems through the theory of Poisson-Nijenhuis manifolds}, Quaderno S/19, Milan, 1984. Re-issued: Universit\`a di Milano-Bicocca, Quaderno 3, 2008.
Available at https://boa.unimib.it/retrieve/e39773b1-610f-35a3-e053-3a05fe0aac26/quaderno3-2008.pdf

%\bibitem{MagriMorosiRagnisco85} Magri, F., Morosi, C., Ragnisco, O.,
%{\it Reduction techniques for infinite-dimensional Hamiltonian systems: some ideas and applications},
%Comm. Math. Phys. {\bf 99} (1985), 115--140. 

%\bibitem{MorosiPizzocchero96} Morosi, C., Pizzocchero, L., {\it $R$-Matrix Theory, Formal Casimirs and the Periodic Toda Lattice}, 
%J. Math. Phys. {\bf 37} (1996), 4484--4513. 

%\bibitem{NDC-D} Nunes da Costa, J.M., Damianou, P.A., {\it Toda systems and exponents of simple Lie groups}, Bull. Sci. Math. {\bf 125} (2001), 49--69.

%\bibitem{Nunes-Marle} Nunes da Costa, J.M., Marle, C.-M., \emph{Reduction of bi-Hamiltonian manifolds and recursion operators}. In: Differential geometry and applications (Brno, 1995), Masaryk Univ., Brno, 1996, pp. 523--538. 

%\bibitem{Okubo90} Okubo, S., {\it Integrability condition and finite-periodic Toda lattice},
%J. Math. Phys. {\bf 31} (1990), 1919--1928.

%\bibitem{Perelomov-book} Perelomov, A.M., {\it Integrable systems of classical mechanics and 
%Lie algebras. Vol. I}, Birkh\"auser Verlag, Basel, 1990.

\bibitem{Pham} Pham, F., {\it Singularities of integrals}, Universitext, Springer, 2010.

\bibitem{SX} Sti\'enon, M., Xu, P., {\it Poisson Quasi-Nijenhuis Manifolds}, Commun. Math. Phys. {\bf 270} (2007), 709--725.

\bibitem{tempestatondoben} Tondo, G., Tempesta, P., {\it Haantjes structures for the Jacobi-Calogero model and the Benenti systems}, SIGMA {\bf 12} (2016), 023, 18 pages.

\bibitem{tempestatondo} Tempesta, P., Tondo, G., {\it Haantjes algebras and diagonalization}, J. Geom. Phys. {\bf 160} (2021), %103868
21 pages.

\bibitem{tempestatondohigher} Tempesta, P., Tondo, G., {\it Higher Haantjes brackets and integrability}, Comm. Math. Phys. {\bf 389} %, N.3, 
(2022), 1647--1671.

\bibitem{tempestatondoclass} Tempesta, P., Tondo, G., {\it Haantjes algebras of classical integrable systems},
Ann. Mat. Pura Appl. {\bf 201} %, N. 1,
(2022), 57--90.

\bibitem{TondoLag} Tondo, G., {\it Haantjes algebras of the Lagrange Top}, 
Theoret.\ and Math.\ Phys.\ {\bf 196} %(3), 
(2018), 1366--1379.

%\bibitem{Tondo1995} Tondo, G., {\it On the integrability of stationary and restricted flows of the KdV hierarchy}, 
%J. Phys. A {\bf 28} (1995), 5097--5115.

%\bibitem{Turiel} Turiel, F.-J., {\it Classification locale d'un couple de formes symplectiques Poisson-compatibles}, 
%C.R. Acad. Sci. Paris S\'er. I Math. {\bf 308} (1989), 575--578.

%\bibitem{Izu} Vaisman, I., {\it The Geometry of Poisson Manifolds}, Birk\"auser Verlag, Basel, 1994.
%
%\bibitem{Complementary} Vaisman, I., {\it Complementary $2$-forms of Poisson structures}, Compositio Math. {\bf 101} (1996), 55--75.

\bibitem{Zucchini} Zucchini, R., {\it The Hitchin model, Poisson-quasi-Nijenhuis, geometry and symmetry reduction}, 
J. High Energy Phys. {\bf 2007}, %no. 10, 075
29 pages.

\end{document}

%%%%%%%%%%%%%%%%%%%%%%%%%
